% One-Loop Health–Education System — White Paper
% Front Matter (LaTeX)
\documentclass[12pt]{article}

% Basic packages
\usepackage[margin=1in]{geometry}
\usepackage{hyperref}
\usepackage{longtable}
\usepackage{booktabs}
\usepackage{enumitem}
\usepackage{setspace}
\usepackage{graphicx}
\usepackage[T1]{fontenc}
\usepackage[utf8]{inputenc}
\usepackage{microtype}
\setstretch{1.15}
\usepackage{multicol}
% ---------- Document Variables ----------
\newcommand{\papertitle}{One-Loop Health–Education System}
\newcommand{\paperversion}{v0.1 \,(Draft)}
\newcommand{\paperdate}{October 24, 2025}
\newcommand{\paperauthors}{Ryan O. van Gelder}
\newcommand{\paperaffiliations}{Citizen Gardens \\ Intrinsica, \& UOR Foundation}
\newcommand{\paperpartners}{Jimmy Danella, Juan Carlos Rodriguez, \\ Kara Olivarria, Gianluca R. Pisano, \& Tyler van Osdol}

% ---------- Title Page ----------
\begin{document}

\begin{titlepage}
  \centering
  {\LARGE \bfseries \papertitle\\[0.75em]}
  {\large White Paper}\\[2em]

  {\large Version: \paperversion}\\
  {\large Date: \paperdate}\\[2em]

  {\normalsize \textbf{Authors}}\\
  {\normalsize \paperauthors}\\[0.75em]

  {\normalsize \textbf{Affiliations}}\\
  {\normalsize \paperaffiliations}\\[0.75em]

  {\normalsize \textbf{Partners}}\\
  {\normalsize \paperpartners}\\[3em]

  % Optional: add a placeholder for partner or program marks
  % \includegraphics[width=0.5\textwidth]{path/to/logo.png}

  \vfill
  {\footnotesize This white paper is provided for planning and discussion. It does not offer medical, legal, or educational advice. See Compliance Notice for details.}
\end{titlepage}

% ---------- Abstract ----------
\begin{abstract}
This white paper describes the One-Loop Health–Education System. The goal is simple: help people learn how to care for themselves while receiving care that learns from them. We bring health and education together on one shared data and AI plane. On this plane we run two clear, coordinated experiences: a \emph{care skin} for patients, families, and clinicians, and a \emph{learning skin} for students, educators, and workforce programs. Both skins use the same trusted data foundation, the same safety rules, and the same feedback cycle.

The system closes the loop from context to outcomes. First, it gathers signals from daily life and care settings with strong privacy controls. Second, it turns those signals into simple, useful insights. Third, it supports actions that people, clinicians, and educators can take. Fourth, it checks what happened and learns from the results. These steps repeat in a steady cycle so the experience improves over time. This closed loop is understandable, auditable, and governed by shared rules across health and learning.

We focus on plain language and safe defaults. Data use is explained in simple terms. People can set limits, change their minds, and receive clear records of how their data moved and why. The design supports local needs while keeping a common core so sites can share what works. The approach is modular, so organizations can adopt only the parts they need and add more over time. The first releases prioritize high-value, low-risk uses such as patient education, adherence support, care transitions, and school-based wellness and literacy programs.

The One-Loop model aims for measurable benefits. For people and families, it aims to improve understanding, confidence, and health habits. For clinicians and care teams, it aims to raise engagement and reduce preventable risk. For students and educators, it aims to increase health literacy and support real-world learning. For communities and systems leaders, it aims to reduce inequities with transparent metrics and fair access to tools. Throughout, the system respects clinical judgment, teacher autonomy, cultural context, and community governance.
\end{abstract}


\newpage
 \begin{multicols}{2}
\begin{singlespace}
\tableofcontents
\end{singlespace}
\end{multicols}
\newpage
% ---------- Compliance Notice and Scope ----------
\section*{Compliance Notice and Scope}
\subsection*{Compliance Notice}
This document explains a design and operating model. It does not provide medical advice, diagnosis, or treatment. It does not provide legal advice or create legal obligations. Any deployment must comply with applicable laws and policies in each jurisdiction and setting. Typical frameworks include, but are not limited to:
\begin{itemize}[leftmargin=*]
  \item Health privacy and security requirements (for example, HIPAA in the United States) for protected health information handled by covered entities and their business associates.
  \item Student privacy and records requirements (for example, FERPA in the United States) for education records and directory information.
  \item Data protection laws that govern consent, purpose limits, individual rights, and international transfers (for example, GDPR in the European Union and similar laws elsewhere).
  \item Standards for information security management and service providers (for example, ISO/IEC 27001 and SOC~2), and organizational controls such as access management, encryption, logging, and incident response.
  \item Clinical safety, human factors, and quality management practices appropriate to the risk level of each use, with human oversight and clear escalation paths.
  \item Accessibility and anti-discrimination requirements to ensure inclusive design and equitable access.
\end{itemize}
All examples in this paper are illustrative. Implementers must complete local risk assessments, data protection impact assessments as needed, and institutional reviews where human subjects or students are involved. Governance structures must include community representation and clear accountability.

\subsection*{Scope}
This white paper covers the purpose, design, and first-phase implementation of a shared health–education loop. It explains the core concepts, the main components, the roles and responsibilities, and the feedback cycle that connects them. It focuses on practical steps that sites can take within 90 days to begin learning safely from real use while protecting people’s rights. It describes the data flows at a high level, the privacy and consent controls, and the measures used to track quality, safety, and equity.

Out of scope are activities that require separate certification or regulatory filings, detailed institution-specific policies, and custom integrations not described in the reference interfaces. Advanced analytics that carry higher risk remain outside the first releases until governance, evaluation, and monitoring capacity are in place. The paper does not set binding policy; it proposes a shared model and a set of operating commitments that organizations can adopt and adapt.

\subsection*{Document Control}
For working drafts, teams may track changes and approvals using a simple table.

\begin{longtable}{@{}p{0.18\linewidth}p{0.18\linewidth}p{0.28\linewidth}p{0.28\linewidth}@{}}
\toprule
\textbf{Version} & \textbf{Date} & \textbf{Change Summary} & \textbf{Authors / Approvers} \\
\midrule
\endhead
\paperversion & \paperdate & Initial draft of front matter and abstract & \paperauthors \\
\bottomrule
\end{longtable}

\newpage
% ---------- Executive Summary ----------
\section{Executive Summary}

\subsection*{The Problem: Fragmented Care and Learning Signals}
People receive care in clinics, hospitals, homes, and community sites. They also learn in schools, training programs, and daily life. Each setting collects pieces of information, but the pieces rarely connect. Health teams may not see what learners understand or struggle with. Educators may not see what health plans ask people to do. Records live in different systems with different rules. Messages arrive too late or not at all. People get duplicate requests, confusing guidance, and low-value alerts. Trust erodes when data moves without clear purpose or when advice is hard to understand or act on. Fragmentation wastes time, hides risk, and widens inequities.

\subsection*{The One-Loop Solution}
The One-Loop Health–Education System ties these signals into a single, governed feedback cycle that is understandable and auditable. The loop follows six repeated steps:
\begin{enumerate}[leftmargin=*]
  \item \textbf{Context intake.} Collect signals from care visits, home use, daily activities, and learning interactions. Use plain consent and strong defaults. Capture what matters and avoid noise.
  \item \textbf{Features and simple risk flags.} Convert raw signals into meaningful features such as adherence trends, symptom changes, and learning progress. Summarize risk in clear language that a person and a professional can both read.
  \item \textbf{Actions.} Deliver small, specific next steps. Examples include micro-lessons, checklists, reminders, care plan tasks, and outreach prompts for teams. Align actions with a person’s goals, language, culture, and access.
  \item \textbf{Outcome check.} Measure whether actions were taken and what changed. Confirm with people, not only with devices or systems. Record both success and friction.
  \item \textbf{Retrain and improve.} Update content, rules, and models using what the outcome check revealed. Keep a record of what changed and why. Roll out improvements safely with staged releases.
  \item \textbf{Equity dashboards.} Track access, engagement, and outcomes by relevant groups approved through governance. Detect gaps early. Guide investments and adjustments.
\end{enumerate}
The loop is simple by design. It reduces guesswork, shortens feedback time, and makes accountability visible. Every step is bounded by consent, purpose, and role-based access.

\subsection*{What Ships First and Why}
Early releases focus on high-value, lower-risk uses that prove the loop and build trust:
\begin{itemize}[leftmargin=*]
  \item \textbf{Patient and family education.} Short, verified lessons linked to a person’s care plan. Readable on low-cost devices. Available offline where needed.
  \item \textbf{Care transitions.} Discharge and handoff checklists that confirm understanding, medications, follow-ups, and warning signs. Outreach prompts for care teams when confirmation fails.
  \item \textbf{Adherence and self-care support.} Simple routines with reminders, reflection questions, and community resources. Respectful nudges rather than spam.
  \item \textbf{School and workforce health literacy.} Modules that teach bio-literacy, mental wellness skills, and safe technology use. De-identified by default and separate from clinical records.
  \item \textbf{Appointment preparation.} Pre-visit questions and expectations so people and teams arrive aligned on goals and decisions.
\end{itemize}
These ship first because they solve widespread pain points, require minimal custom integration, and demonstrate measurable gains in understanding, engagement, and safety. They also keep protected health information within clinical boundaries while learning spaces work with de-identified or aggregate data. Sites can start small, run weekly reviews, and scale what works.

\subsection*{Service Level Objective (SLO) Targets}
Initial targets guide operations and set clear expectations. They will be refined with each site:
\begin{itemize}[leftmargin=*]
  \item \textbf{Availability.} The core services remain available during agreed service hours, with planned maintenance windows announced in advance.
  \item \textbf{Action delivery time.} Most actions, lessons, or prompts appear within a few seconds after the triggering event is received.
  \item \textbf{Data freshness.} Updates from connected clinical systems and learning platforms appear in near real time or within minutes, depending on partner systems.
  \item \textbf{Audit completeness.} Every inbound signal, decision, action, and change is logged with who, what, when, and purpose of use.
  \item \textbf{Explanation coverage.} Every insight or recommendation includes a plain-language “why this, why now” note and a link to supporting sources the user is allowed to see.
  \item \textbf{Issue response.} Confirmed safety or privacy issues are triaged immediately, with user notification and mitigation steps defined in the runbook.
\end{itemize}
These targets keep the loop responsive and trustworthy without over-promising. They are non-binding for design drafts and will be formalized during site onboarding.

\subsection*{Privacy and Safety Guardrails}
Guardrails shape what the system will and will not do:
\begin{itemize}[leftmargin=*]
  \item \textbf{Purpose limits and consent.} Each use is tied to a clear purpose. People can see, change, or revoke consent. Receipts confirm how data moved and why.
  \item \textbf{Boundary control.} Clinical spaces keep protected health information under clinical governance. Learning spaces receive de-identified or aggregate data unless a lawful and explicit exception is approved.
  \item \textbf{Least data, shortest time.} Collect the minimum necessary data and keep it only as long as needed for the stated purpose.
  \item \textbf{Access controls.} Role-based and attribute-based access limit who can see what. Administrative access requires strong authentication and is time-boxed.
  \item \textbf{Human oversight.} Clinicians and educators remain the decision makers. The system proposes actions and explains them; humans accept, modify, or decline.
  \item \textbf{Safety reviews and kill switch.} New or changed content and models pass review before broad release. A kill switch disables any component that causes harm or operates outside policy.
  \item \textbf{Fairness checks.} Equity dashboards monitor access, engagement, and outcomes. If gaps appear, the system slows or stops rollouts until mitigations are in place.
  \item \textbf{Transparency.} People can read plain-language privacy notices, see data sources allowed for their experience, and download summaries of their interactions.
\end{itemize}
These guardrails are built into daily operations, not added later. They make the loop safe to use at small and large scale.

\bigskip
% ---------- Purpose and Operating Tenets ----------
\section{Purpose and Operating Tenets}

\subsection*{Purpose}
The purpose is to close the loop between clinical care and learning under shared governance. The system treats health and education as one coordinated service that listens, acts, and learns together. It connects day-to-day signals, care plans, and learning activities so that people receive timely support and so teams see what works. Decisions about data, models, and content are made with the people who are affected by them. The aim is to achieve better understanding, safer care, and fairer outcomes while keeping authority with clinicians, educators, and communities.

\subsection*{Operating Tenets}
\begin{enumerate}[leftmargin=*]
  \item \textbf{Simple, purposeful use.} Every data use has a clear purpose that benefits the person and the site. If a purpose is unclear, the system does not use the data.
  \item \textbf{Human oversight.} Clinicians and educators remain in charge. The system explains its suggestions in plain language and offers easy ways to accept, adjust, or decline.
  \item \textbf{Least data, least friction.} Collect and retain only what is needed. Ask once, reuse responsibly, and avoid duplicate requests or alerts.
  \item \textbf{Interoperable by default.} Use open formats and clear contracts so partners can connect without custom one-off work.
  \item \textbf{Equity at the core.} Monitor access, engagement, and outcomes. Slow or stop rollouts if harm or gaps appear, and invest where support is needed most.
  \item \textbf{Plain language and transparency.} Policies, explanations, and receipts are readable by non-experts. People can see how their data moved and why.
  \item \textbf{Local control with shared standards.} Sites keep governance and control over their environments while adopting a common model that allows learning across sites.
\end{enumerate}

\subsection*{Governance and Change Control}
Governance combines clear roles, routine reviews, and reliable change processes:
\begin{itemize}[leftmargin=*]
  \item \textbf{Steering Council.} Sets priorities, approves major uses, and validates measures of benefit and risk. Includes clinical, education, community, and operations leaders.
  \item \textbf{Data and Safety Board.} Oversees privacy, security, safety, and evaluation. Reviews incidents and approves model or content releases that affect risk.
  \item \textbf{Community Advisory Circle.} Provides lived-experience perspective, language access guidance, and cultural relevance checks. Reviews communications and consent flows.
  \item \textbf{Site Working Group.} Runs pilots, manages integrations, and tracks day-to-day issues. Publishes a shared runbook and weekly change log.
\end{itemize}

\paragraph{Change Classes}
\begin{itemize}[leftmargin=*]
  \item \textbf{Standard changes.} Low-risk, pre-approved steps such as content tweaks or display fixes. Batched and announced.
  \item \textbf{Minor changes.} Limited risk and scope, such as a new checklist or a revised dashboard tile. Require peer review and staged rollout.
  \item \textbf{Major changes.} Affect safety, privacy, or core logic. Require council approval, sandbox testing, impact assessment, and explicit rollback plans.
\end{itemize}

\paragraph{Change Process}
Propose → review → sandbox trial → staged rollout → monitor → decide to keep, adjust, or roll back. Each step has an owner, a checklist, and a timestamped record.

\subsection*{RACI at a Glance}
The matrix below clarifies who is Responsible (R), Accountable (A), Consulted (C), and Informed (I) for common activities. Roles are examples and can be mapped to local titles.

\begin{longtable}{@{}p{0.34\linewidth}p{0.12\linewidth}p{0.12\linewidth}p{0.12\linewidth}p{0.12\linewidth}@{}}
\toprule
\textbf{Activity} & \textbf{Clinical Lead} & \textbf{Education Lead} & \textbf{Data Steward} & \textbf{Security Officer} \\
\midrule
\endhead
Approve new use of data & C & C & A & C \\
Change learning content & I & A & C & I \\
Adjust data sources or mappings & C & I & A & C \\
Incident response and user notice & C & C & R & A \\
Release new model or rule set & C & C & R & A \\
Decommission a feature & R & R & C & C \\
\bottomrule
\end{longtable}

Additional roles often included: Community Representative (C/I), Product Owner (A for backlog), Site Admin (R for configuration), and Vendor PM (I for cadence and reporting).

\subsection*{Privacy-by-Design and DID-Bound Consent}
Consent and privacy controls are embedded in the workflow, not bolted on later:
\begin{itemize}[leftmargin=*]
  \item \textbf{DID-bound consent.} Each person can have a privacy identity anchored to a decentralized identifier (DID). Consents are linked to that identifier and to a clear purpose, so they travel with the person across participating systems.
  \item \textbf{Layered consent.} Present short, plain choices with links to deeper detail. Offer separate switches for data types and uses. Make revocation easy and honored quickly.
  \item \textbf{Receipts for data use.} After important actions, issue a readable receipt that shows what data moved, for what purpose, and who saw it. Store copies for audits.
  \item \textbf{Data minimization.} Prefer de-identified or aggregate data for learning spaces. Keep protected health information under clinical governance unless a lawful, explicit exception applies.
  \item \textbf{Clear boundaries.} Education platforms and clinical systems exchange only what is necessary through well-defined interfaces with role-based access.
\end{itemize}

\subsection*{Deterministic Traceability}
Every event leaves a trail that people and auditors can follow. Inbound signals, transformations, decisions, actions, and outcomes are logged with who did what, when, and for which purpose. The same identifier follows the event across services so teams can reconstruct a timeline without guesswork. Regular reviews check that records are complete and that explanations remain accurate and readable.

\subsection*{BYOC Visit Pack}
The Bring-Your-Own-Context (BYOC) Visit Pack is a small, portable bundle that travels with a person before, during, and after a visit:
\begin{itemize}[leftmargin=*]
  \item \textbf{Before the visit.} A short briefing explains what to expect, what to bring, and the decisions likely to be made. People can add goals, questions, accessibility needs, and language preferences.
  \item \textbf{During the visit.} The pack provides checklists, visuals, and translation aids. It captures agreed actions and notes in simple terms that the person approves before saving.
  \item \textbf{After the visit.} The pack delivers follow-up steps, micro-lessons tied to the plan, and clear ways to ask for help. It can be shared with trusted caregivers or educators using consent the person controls.
  \item \textbf{Portability and privacy.} The pack works on low-cost phones and public computers. It uses short-lived links or codes and stores only the minimum needed on devices.
\end{itemize}

\subsection*{Offline-First Sovereignty}
Many people and sites have limited or unreliable connectivity. The system is usable offline and syncs safely when a connection returns:
\begin{itemize}[leftmargin=*]
  \item Core features work without a network, including viewing key materials, capturing notes, and checking off tasks.
  \item Data is queued locally and encrypted. Sync respects consent and purpose limits.
  \item Conflicts are handled with clear prompts and side-by-side comparisons. Users choose what to keep.
  \item Sites can host services locally or in a preferred region to meet policy or community requirements.
\end{itemize}

\subsection*{Safety and Ethics}
Safety and ethics guide daily operations and releases:
\begin{itemize}[leftmargin=*]
  \item Define what harm means for this context and track leading indicators, not only outcomes.
  \item Require human review for sensitive content and for any automated recommendation that could change care or learning trajectory.
  \item Test content and features with diverse users, including those with disabilities and limited digital access.
  \item Run fairness checks and document mitigations before broad release. Pause when results are uncertain.
  \item Maintain a public-facing change log and a private incident log with timely user notices when appropriate.
\end{itemize}

% End Purpose and Operating Tenets

% ---------- System Synopsis ----------
\section{System Synopsis}

\subsection*{Roles at a Glance}
This section names the core parts of the system and how they work together. Each part has a clear job, clear inputs and outputs, and clear owners.

\begin{longtable}{@{}p{0.18\linewidth}p{0.34\linewidth}p{0.18\linewidth}p{0.24\linewidth}@{}}
\toprule
\textbf{Component} & \textbf{Core Role} & \textbf{Primary Users} & \textbf{Key Inputs \& Outputs} \\
\midrule
\endhead
Intrinsica (Data/AI) &
Shared data and intelligence plane. Turns raw signals into usable insights and simple risk flags. Routes events, keeps feature stores, manages model versions, and records explanations. &
Data stewards, analysts, integration engineers, product teams. &
\emph{Inputs:} de-identified learning telemetry, clinical events, device and app signals, consent states. \newline
\emph{Outputs:} insights, triggers, risk flags, audit records, explanation notes, de-identified research exports. \\
\addlinespace
S2H (Clinic Apps) &
Point-of-care and follow-up apps for patients and care teams. Delivers checklists, reminders, and care plan tasks. Captures confirmations and barriers. Works online and offline. &
Clinicians, care coordinators, patients, caregivers. &
\emph{Inputs:} care plans, triggers from Intrinsica, BYOC Visit Pack entries. \newline
\emph{Outputs:} task status, symptom updates, adherence notes, escalation prompts. \\
\addlinespace
Hologram (Record) &
Context-first longitudinal record. Holds visit summaries, plans, consents, and person-stated goals. Serves as the clinical source of truth and boundary for protected information. &
Clinical teams, health information management, compliance. &
\emph{Inputs:} EHR events, S2H confirmations, signed consents. \newline
\emph{Outputs:} updated plans, care summaries, consent receipts, role-limited views for partners. \\
\addlinespace
HEP/ULMS (Education Layer) &
Learning experience and content services. Delivers short, verified lessons and practice activities linked to a person’s context. Collects de-identified progress data for improvement and equity tracking. &
Students, patients, families, educators, workforce programs. &
\emph{Inputs:} learning goals, triggers from Intrinsica, language and accessibility preferences. \newline
\emph{Outputs:} lesson completions, feedback, engagement metrics, de-identified cohort reports. \\
\addlinespace
ADEC–SCM (System Coordinator) &
Operational coordination and guardrails. Maintains system state, applies policy checks, sequences actions, and pauses or slows features when safety or fairness signals degrade. &
Site operators, safety officers, governance leads. &
\emph{Inputs:} policy rules, live service signals, exception reports. \newline
\emph{Outputs:} allow/deny decisions, rollout gates, incident hooks, stability notices. \\
\addlinespace
MQEM/PIRTM (Quality \& Provenance Engine) &
Quality, drift, and provenance services. Cleans noisy data, checks model behavior over time, traces data lineage, and controls staged releases with audit-ready records. &
Data quality teams, model owners, auditors. &
\emph{Inputs:} raw and processed data samples, model metrics, release candidates. \newline
\emph{Outputs:} cleaned datasets, drift alerts, lineage graphs, release approvals or blocks. \\
\bottomrule
\end{longtable}

\subsection*{One-Loop Concept Diagram (Text)}
The diagram below shows how information moves. It highlights what is read and what is written at each step, and where consent and privacy rules apply.

\begin{verbatim}
[People, Families, Teams, Schools, Community]
           |
           | 1) Signals and actions in daily life and care
           v
+-------------------+        Read: context, plans, consents
|   Hologram Record |------> Write: visit summaries, plan updates
+-------------------+        Boundary for protected health info
           |
           | Trusted, governed events (no unnecessary identifiers)
           v
+-------------------+        Read: events, progress, preferences
|  Intrinsica Data  |------> Write: insights, triggers, explanations
|     and AI Plane  |        (kept with audit trails)
+-------------------+
     |           |
     |           +-------------------------------+
     |                                           |
     v                                           v
+-------------------+                    +--------------------+
|   S2H Clinic Apps |  Read: triggers    |  HEP/ULMS Learning |
|  (patient & team) |  Write: task status|   (lessons & LXP)  |
+-------------------+       and notes    +--------------------+
     |                                           |
     |  3) Confirm actions and barriers          |  4) Capture de-identified learning data
     |                                           |
     +--------------------+  --------------------+
                          v  v
                 +--------------------+
                 |  Equity Dashboards |
                 |  & Governance View |
                 +--------------------+
                          |
          5) Review gaps, change content, adjust rollout
                          |
                          v
                 +--------------------+
                 | MQEM/PIRTM Quality |
                 |  & Provenance Gate |
                 +--------------------+
                          |
          6) Safe releases, versioned and reversible
                          |
                          v
                 +--------------------+
                 |  ADEC–SCM System   |
                 |   Coordinator      |
                 +--------------------+
                          |
          7) Policy checks, stability signals, kill switch
                          |
                          v
                 +--------------------+
                 |  Intrinsica Data   |
                 |   and AI Plane     |
                 +--------------------+
                          |
          8) Updated insights and triggers flow back
                          |
                          v
        [People, Families, Teams, Schools, Community]
\end{verbatim}

\paragraph{How the pieces align}
\begin{itemize}[leftmargin=*]
  \item Hologram holds the clinical source of truth and consent records. It is the entry and exit point for protected information.
  \item Intrinsica converts events into simple, useful signals and dispatches triggers with clear explanations and audit entries.
  \item S2H apps turn triggers into checklists and support. They capture confirmations and barriers in ways that are easy for people to use.
  \item HEP/ULMS delivers short, accessible learning tied to real decisions. It returns only de-identified progress data unless an approved exception applies.
  \item ADEC–SCM coordinates the system and enforces safety, fairness, and pacing rules across sites.
  \item MQEM/PIRTM keeps the data trustworthy, the models stable, and every release traceable and reversible.
\end{itemize}

% End System Synopsis


% ---------- Architecture Overview ----------
\section{Architecture Overview}

\subsection*{Layer Index}
The system is organized into five layers. Each layer has a focused job and clear boundaries.

\begin{itemize}[leftmargin=*]
  \item \textbf{Experience layer} — Interfaces for people, families, clinicians, educators, and operators. Web, mobile, kiosk, and offline packs.
  \item \textbf{Services layer} — Application services that turn insights into actions. Orchestration, content, notifications, tasking, and help.
  \item \textbf{Data/AI layer} — Trusted data foundation. Storage, feature pipelines, simple risk flags, explanations, quality checks.
  \item \textbf{Interop layer} — Connectors and contracts. Clinical and learning standards, webhooks, queues, routing, and de-identification.
  \item \textbf{Governance layer} — Policy, consent, audit, release control, evaluation, and equity dashboards.
\end{itemize}

\subsection*{Experience Layer}
\textbf{Purpose.} Deliver clear, accessible experiences that help people act and understand.  
\textbf{Scope.} Patient and family apps, team tools, educator portals, and community touchpoints. Works on low-cost devices and in low-connectivity settings.  
\textbf{Responsibilities.}
\begin{itemize}[leftmargin=*]
  \item Present actions, checklists, and short lessons at the right time.
  \item Capture confirmations, barriers, preferences, and feedback.
  \item Support BYOC Visit Packs for pre-visit, in-visit, and post-visit use.
  \item Provide readable consent choices, notices, and receipts.
  \item Operate offline when needed and sync safely later.
\end{itemize}
\textbf{Inputs.} Triggers, plans, learning assignments, and explanations from upstream services.  
\textbf{Outputs.} Task status, notes, lesson completions, and support requests.

\subsection*{Services Layer}
\textbf{Purpose.} Translate insights into practical steps and coordinate delivery across channels.  
\textbf{Capabilities.}
\begin{itemize}[leftmargin=*]
  \item \textit{Action orchestration} to create tasks, checklists, messages, and micro-lessons.
  \item \textit{Content services} to store versioned lessons, visuals, and checklists with accessibility metadata.
  \item \textit{Notification and scheduling} for reminders, follow-ups, and handoffs.
  \item \textit{Consent and identity UX} for readable choices and secure sharing with trusted supporters.
  \item \textit{Support desk} workflows for questions, issues, and escalation.
\end{itemize}
\textbf{Inputs.} Insights and triggers from Data/AI, policies from Governance.  
\textbf{Outputs.} Action payloads to Experience, acknowledgments and outcomes back to Data/AI.

\subsection*{Data/AI Layer}
\textbf{Purpose.} Turn raw signals into clean data, simple features, and understandable insights with full traceability.  
\textbf{Capabilities.}
\begin{itemize}[leftmargin=*]
  \item Event and record stores with clear retention and purpose tags.
  \item Feature pipelines that summarize trends such as adherence or engagement.
  \item Simple risk flags with plain-language explanations and links to sources the user can see.
  \item Quality, drift, and lineage checks to keep outputs stable and reproducible.
  \item Safe exports for de-identified research and aggregated learning reports when approved.
\end{itemize}
\textbf{Inputs.} Normalized events from Interop and outcomes from Services and Experience.  
\textbf{Outputs.} Insights, triggers, explanations, and audit records.

\subsection*{Interop Layer}
\textbf{Purpose.} Connect partner systems and keep boundaries clear.  
\textbf{Capabilities.}
\begin{itemize}[leftmargin=*]
  \item Standards-based connectors for clinical and learning platforms.
  \item Webhooks and queues for near real-time event flow with retries and dead-letter handling.
  \item De-identification and field-level filtering so only necessary data crosses boundaries.
  \item Throttling and load-shedding to protect partner systems during peaks.
  \item Schema and contract registry so partners can self-serve integrations.
\end{itemize}
\textbf{Inputs.} Events, records, and consent states from external systems.  
\textbf{Outputs.} Normalized, purpose-limited events and acknowledgments to partners.

\subsection*{Governance Layer}
\textbf{Purpose.} Keep the system lawful, fair, and accountable.  
\textbf{Capabilities.}
\begin{itemize}[leftmargin=*]
  \item Policy rules that gate features and releases.
  \item Consent records, access reviews, and audit trails that people and auditors can read.
  \item Evaluation plans, safety checks, and rollback controls.
  \item Equity dashboards that surface gaps in access, engagement, and outcomes.
  \item Transparent change logs and incident workflows with timely user notices when appropriate.
\end{itemize}
\textbf{Inputs.} Measures, incidents, and proposed changes from all layers.  
\textbf{Outputs.} Approvals, pauses, rollbacks, and public reports.

\subsection*{End-to-End Flow: Signals to Dashboards}
This narrative shows how information moves through the layers in everyday use.

\begin{enumerate}[leftmargin=*]
  \item \textbf{Signals arrive.} A person completes a short lesson, checks off a task, or reports a symptom. A clinic system emits a discharge event. A school system records a wellness activity.
  \item \textbf{Interop normalizes.} Connectors receive the events, apply field filters, remove unnecessary identifiers, and attach purpose tags. The events enter the shared stream.
  \item \textbf{Data/AI ingests.} The layer validates formats, stores the events, and updates simple features such as recent activity, comprehension indicators, or follow-up status.
  \item \textbf{Insights are created.} Rules generate clear, human-readable insights and triggers. Each has an explanation note and a link to related materials.
  \item \textbf{Services orchestrate.} The orchestration service selects the best next action: a checklist item, a reminder, a short lesson, or a team outreach prompt. It schedules delivery and sets a soft due date.
  \item \textbf{Experience delivers.} People see the action in their app or receive it by the chosen channel. The interface shows why it matters and how long it takes.
  \item \textbf{Outcomes return.} The app or team tool captures completion, barriers, or requests for help. These outcomes flow back through Services to Data/AI.
  \item \textbf{Learning and improvement.} Data/AI updates features and adjusts thresholds or content choices based on what worked. Proposed changes move to Governance for review when they affect wider groups.
  \item \textbf{Dashboards update.} Governance receives refreshed measures. Equity dashboards highlight any emerging gaps by cohort approved through policy. Operators see stability and quality indicators.
  \item \textbf{Release control.} If results look good, Governance allows wider rollout. If issues appear, it slows or pauses the change and opens an incident or improvement task.
\end{enumerate}

\subsection*{Separation of Concerns and Boundaries}
\begin{itemize}[leftmargin=*]
  \item Protected health information remains within clinical boundaries. Learning spaces work with de-identified or aggregated data unless a specific, lawful exception is approved.
  \item Each layer owns its logs and exports summaries the next layer needs. No layer reads more than it must.
  \item Every action and change includes a human-readable explanation and a record of who initiated it and for what purpose.
\end{itemize}

\subsection*{Observability and Operations}
\begin{itemize}[leftmargin=*]
  \item Unified logging across layers with trace identifiers so teams can follow an event from intake to outcome.
  \item Live dashboards for availability, delivery time, data freshness, audit completeness, and explanation coverage.
  \item Health checks, circuit breakers, and a kill switch for any component that misbehaves.
\end{itemize}

% End Architecture Overview

% ---------- Intrinsica Data/AI Plane ----------
\section{Intrinsica Data/AI Plane}

\subsection*{Role and Scope}
Intrinsica is the shared data and intelligence plane. It turns raw events into clear insights and safe triggers, then records why each insight was produced. It anchors auditability, quality, and release control for the full loop.

\subsection*{Pipeline Overview}
The pipeline follows a simple sequence that repeats continuously:
\begin{enumerate}[leftmargin=*]
  \item \textbf{Ingest} — receive purpose-tagged events and records.
  \item \textbf{Feature store} — create and update simple, reusable summaries.
  \item \textbf{Risk registry} — register plain-language risk flags with explanations.
  \item \textbf{Digital twin and simulation} — stage safe dry-runs of new content and rules against recent history before any broad release.
\end{enumerate}
Each stage writes a complete audit entry that names the input sources used, the purpose, and the person or role that initiated the step.

\subsection*{Ingest}
\begin{itemize}[leftmargin=*]
  \item Accepts normalized events from the interop layer with purpose tags and minimal fields.
  \item Validates format, consent status, and boundary rules before storing.
  \item Assigns a trace identifier so every event can be followed end to end.
  \item Queues are durable with retry, backoff, and dead-letter handling to avoid data loss.
\end{itemize}

\subsection*{Feature Store}
\begin{itemize}[leftmargin=*]
  \item Maintains simple, human-meaningful summaries such as recent activity, adherence streaks, comprehension indicators, and follow-up status.
  \item Features are versioned. Each has a short description, owner, inputs, and refresh policy.
  \item Only approved features are available to downstream services. Draft features stay isolated until reviewed.
\end{itemize}

\subsection*{Risk Registry}
\begin{itemize}[leftmargin=*]
  \item Holds current risk flags with a clear title, level, rationale, and suggested next step.
  \item Every flag includes a readable “why this, why now” explanation and links to sources the user is allowed to see.
  \item Flags expire automatically unless reaffirmed by fresh signals or a human check.
  \item Changes to thresholds or rules are staged, reviewed, and rolled out gradually.
\end{itemize}

\subsection*{Digital Twin and Simulation}
\begin{itemize}[leftmargin=*]
  \item Mirrors live data flows in an isolated space using recent, de-identified samples.
  \item Runs proposed content, rules, or model releases as dry-runs to estimate user impact, delivery volume, and potential equity effects.
  \item Produces a short release note with predicted outcomes and any guardrails required before approval.
\end{itemize}

\subsection*{Interop and Event Bus}
\begin{itemize}[leftmargin=*]
  \item Subscribes to a shared event bus that carries clinical, learning, and user-generated events with purpose tags.
  \item Publishes insights, triggers, and acknowledgments back to the bus for consumption by services and apps.
  \item Enforces rate limits and back-pressure so partner systems are protected during peak loads.
  \item Uses idempotent message keys so retries never create duplicates.
\end{itemize}

\subsection*{De-Identification and Exports}
\begin{itemize}[leftmargin=*]
  \item Applies field-level filtering and tokenization before any data crosses the clinical boundary.
  \item Provides safe, de-identified exports for research and learning reports with documented risks and residuals.
  \item All exports are logged with purpose, requester, fields included, and retention period.
\end{itemize}

\subsection*{Trust Fabric}
\begin{itemize}[leftmargin=*]
  \item \textbf{Decentralized identifiers (DIDs).} Bind consent and sharing choices to a portable identifier under user or site control.
  \item \textbf{Purpose-bound consent.} Every read and write cites an allowed purpose. Revocation removes access for future processing.
  \item \textbf{Ledgered provenance.} Key decisions and releases produce tamper-evident receipts that include inputs, owners, and timestamps.
  \item \textbf{Sovereign sync.} Sites can run local nodes and sync only the data and artifacts they approve, on a schedule that fits local policy.
\end{itemize}

\subsection*{Observability and Controls}
\begin{itemize}[leftmargin=*]
  \item Unified logs with trace identifiers from ingest to action delivery.
  \item Real-time views for data freshness, delivery success, backlog depth, and error classes.
  \item Kill switch for any feature that misbehaves or exceeds safety budgets.
  \item One-click rollback to a prior, signed release of features, rules, or content.
\end{itemize}

\subsection*{Data Management}
\begin{itemize}[leftmargin=*]
  \item Clear retention by data type and purpose, with automatic deletion at end of life.
  \item Access is role-based and attribute-based. Administrative access is time-limited and requires step-up authentication.
  \item All datasets and features list owners, stewards, and review dates.
\end{itemize}

\subsection*{Initial SLO Targets}
Targets guide operations and will be refined with each site:
\begin{longtable}{@{}p{0.28\linewidth}p{0.64\linewidth}@{}}
\toprule
\textbf{Area} & \textbf{Target and Notes} \\
\midrule
\endhead
Insight latency & Typical insights and triggers are produced within seconds after an event arrives, subject to partner system delays. \\
Commit durability & Once acknowledged, events and derived updates are durable and visible across services with no data loss on retry. \\
Data freshness & Near real-time updates for active users; batch windows clearly posted for low-priority feeds. \\
Privacy pass-rate & Processing that cites a purpose and consent state succeeds without manual review in the vast majority of cases; exceptions are queued with clear reasons. \\
Audit completeness & Every event and decision carries a trace identifier and a readable explanation note. Missing fields trigger alerts. \\
Issue response & Confirmed safety or privacy issues are triaged immediately with user notice and mitigation steps from the runbook. \\
\bottomrule
\end{longtable}

\subsection*{Release and Change Management}
\begin{itemize}[leftmargin=*]
  \item Changes to features, thresholds, or content start in the digital twin, produce a short release note, and move through review.
  \item Staged rollout with automatic pause if error rates, fairness checks, or user feedback cross agreed thresholds.
  \item Public change log for non-sensitive updates and a detailed internal log for audits.
\end{itemize}

% End Intrinsica Data/AI Plane

% ---------- Hologram Health: Context-First Record ----------
\section{Hologram Health: Context-First Record}

\subsection*{Role and Scope}
Hologram Health is the context-first record for clinical work. It organizes what is known about a person into a clear, living timeline that teams can trust. It holds visit summaries, plans, consents, and person-stated goals. It gives each encounter a short, readable story, then connects that story to next steps. Hologram is the boundary for protected clinical information and the place where people confirm that information is correct.

\subsection*{Agents}
\paragraph{HoloGraph}
HoloGraph is the organizing agent. It links people, visits, plans, medications, documents, and goals into a single view. It resolves duplicates, applies naming standards, and surfaces context cards for quick decisions. When a team opens a chart, HoloGraph selects the most relevant facts for that moment and shows supporting details on demand. It keeps a full change history and makes merges and corrections transparent.

\paragraph{Spectra}
Spectra is the retrieval and explanation agent. It answers plain questions about the record, shows where each fact came from, and offers short, source-linked summaries. Spectra respects role-based views and only returns information the user is allowed to see. It supports voice and text, works on low-cost devices, and degrades gracefully when connectivity is limited.

\subsection*{MVP Experiences}
\begin{itemize}[leftmargin=*]
  \item \textbf{SafeCare.} Safety-critical essentials at a glance: allergies, alerts, active medications, recent changes, and urgent tasks. Comes with short checklists for high-risk moments such as medication reconciliation and discharge.
  \item \textbf{BridgeCare.} Smooth handoffs across settings. Teams confirm contacts, follow-ups, shared plans, equipment needs, and warning signs. BridgeCare produces a simple, shared summary for the person and for receiving teams.
  \item \textbf{PrecisionCare.} Person-centered plan builder. Teams record goals, preferences, and constraints, then assemble a focused plan with small, trackable steps. PrecisionCare links each step to clear explanations and resources.
\end{itemize}

\subsection*{Outputs}
Hologram produces artifacts that humans can use and systems can process.
\begin{itemize}[leftmargin=*]
  \item \textbf{Visit summaries} that are short, structured, and readable.
  \item \textbf{Plan updates} with clear owners, due windows, and review dates.
  \item \textbf{Consent receipts} that show what was agreed, for what purpose, and how to change it later.
  \item \textbf{Care packets} for handoffs, including checklists, instructions, and contact points.
  \item \textbf{Machine-readable messages} that announce key changes to downstream services along with explanation notes and trace identifiers.
\end{itemize}

\subsection*{Interfaces}
\textbf{Clinical workspace.} A focused view for teams to review context, confirm details, and finalize plans.  
\textbf{Patient and caregiver portal.} A simple view to check summaries, ask questions, and share approved information with trusted supporters.  
\textbf{Admin console.} Tools for data stewardship, consent management, and release control.  
\textbf{System hooks.} Standards-based read and write endpoints for partner systems, plus webhooks and queues for near real-time events.  
\textbf{Export tools.} Role-limited exports for referrals, transitions, and de-identified reporting when allowed.

\subsection*{Operating Constraints}
\begin{itemize}[leftmargin=*]
  \item \textbf{Boundary control.} Protected clinical information remains under clinical governance with strict access reviews.
  \item \textbf{Role-limited views.} Users see only what they need for their role and purpose.
  \item \textbf{Accessibility.} Interfaces meet common accessibility standards and support multiple languages.
  \item \textbf{Performance.} Records and summaries open quickly, even on low-cost devices and during peak hours.
  \item \textbf{Break-glass.} Emergency access is available with immediate logging, after-action review, and user notice when appropriate.
  \item \textbf{Versioning.} Every summary, plan, and consent is versioned with who changed what and when.
  \item \textbf{Retention.} Data types carry clear retention periods and automatic end-of-life handling.
\end{itemize}

\subsection*{Success Criteria}
\begin{itemize}[leftmargin=*]
  \item \textbf{Findability.} Time to locate key facts (allergies, medications, follow-ups) meets target thresholds.
  \item \textbf{Clarity.} Summaries score at a plain-language reading level and pass spot checks for accuracy.
  \item \textbf{Completeness.} Required fields for handoffs are present before transitions proceed.
  \item \textbf{Quality.} Fewer duplicate records and fewer conflicting plans over time.
  \item \textbf{Timeliness.} Discharge packets and follow-up details are ready before the person leaves the setting.
  \item \textbf{Trust.} Users report that explanations and receipts are understandable and useful.
  \item \textbf{Equity.} Use and outcomes are balanced across groups defined in governance, with early detection and response to gaps.
\end{itemize}

\subsection*{Clinic-Ready BYOC Flows and Verification}
The Bring-Your-Own-Context (BYOC) flow lets people carry their own context into and out of visits in a safe, simple way.

\paragraph{Pre-visit}
People receive a short invite by their preferred channel with a one-time code or link. They can add goals, questions, accessibility needs, medications, and support contacts. A clear notice explains what will be shared during the visit and how to change or revoke consent.

\paragraph{Check-in}
At arrival, a kiosk, tablet, or staff member verifies the one-time code or a QR on the invite. If the person prefers, they can skip devices and confirm verbally with staff. The system shows a quick summary of what will be shared for the visit and asks for confirmation.

\paragraph{During the visit}
The BYOC pack shows checklists and visuals that match the purpose of the visit. Teams update medications, capture decisions, and mark follow-ups. People can request a pause or decline any item that feels out of scope.

\paragraph{After the visit}
The pack delivers a clear summary and next steps. People can share the summary with trusted supporters or schools using time-limited links tied to consent. If something is wrong, they can flag it for review.

\paragraph{Verification and safety}
\begin{itemize}[leftmargin=*]
  \item One-time codes and short-lived links reduce risk. Codes can be sent by SMS, email, or printed at check-in.
  \item Staff can verify identity using standard clinic procedures when digital verification is not possible.
  \item Caregivers can be added with consent, with simple ways to remove access later.
  \item All confirmations and changes produce receipts and appear in the audit log.
\end{itemize}

% End Hologram Health: Context-First Record

\section{HEP/ULMS Learning Layer}

\subsection*{Role and Scope}
The HEP/ULMS layer delivers short, verified learning experiences that match a person’s real context. It accepts triggers from upstream services, chooses the right lesson or practice activity, and records de-identified progress for improvement and equity tracking. It treats learning as part of daily care, not a separate chore. Content is short, readable, accessible, multilingual, and usable on low-cost devices. Every item carries a clear “why this, why now” note and a time estimate.

\subsection*{Function Overview}
\begin{itemize}[leftmargin=*]
  \item Selects a lesson or practice activity that fits the person’s goals, language, and access.
  \item Delivers the activity in small steps, with optional audio, captions, and images.
  \item Captures completion, comprehension checks, and user feedback in de-identified form.
  \item Returns summarized outcomes to the Data/AI plane and updates dashboards for governance.
  \item Allows local programs to author, translate, and version content with clear review and release steps.
\end{itemize}

\subsection*{Schema Primitives}
The learning layer uses a few simple, reusable primitives. These keep interfaces understandable and make it easy to plug new content or partners into the loop.

\paragraph{PER — Personalized Education Record}
A light record of a person’s learning context and preferences that the site is allowed to use.
\begin{itemize}[leftmargin=*]
  \item \textbf{Purpose.} Keep only what is needed to tailor delivery and accessibility.
  \item \textbf{Typical fields.} Preferred language, reading aid settings, audio preference, device constraints, time-of-day preference, cultural notes, accessibility needs, consent scope for learning use.
  \item \textbf{Produced by.} Intake flows, BYOC packs, and staff edits with consent.
  \item \textbf{Consumed by.} HEP/ULMS delivery and orchestration services.
\end{itemize}

\paragraph{MLEM — Modular Learning Element Model}
The unit of learning that can be mixed and matched across programs.
\begin{itemize}[leftmargin=*]
  \item \textbf{Purpose.} Package a short lesson or practice step with clear outcomes and media.
  \item \textbf{Typical fields.} Title, objective in plain language, estimated time, accessibility notes, media assets, comprehension check, version, owner, required CSL.
  \item \textbf{Produced by.} Content teams with review and localization steps.
  \item \textbf{Consumed by.} HEP/ULMS runtime, S2H apps, and summaries.
\end{itemize}

\paragraph{TCI — Task–Content Interface}
The contract that links a real-world task to a learning element.
\begin{itemize}[leftmargin=*]
  \item \textbf{Purpose.} Ensure that every lesson is tied to an action people may take.
  \item \textbf{Typical fields.} Task name, success criteria, helpful materials, step-by-step guide, escalation route, related MLEM IDs.
  \item \textbf{Produced by.} Program designers with clinical or educator review.
  \item \textbf{Consumed by.} Orchestration services and team consoles.
\end{itemize}

\paragraph{CSLState — Context and Safety Level State}
The gating state that controls when to deliver a recommendation.
\begin{itemize}[leftmargin=*]
  \item \textbf{Purpose.} Deliver only when context fit, safety, and explanation quality meet the local threshold.
  \item \textbf{Typical fields.} Level (\textit{green/amber/red}), reason, prerequisites to lift a block, allowed channels, next review time.
  \item \textbf{Produced by.} Intrinsica rules and operator settings.
  \item \textbf{Consumed by.} HEP/ULMS runtime and S2H action selection.
\end{itemize}

\paragraph{ReciprocityDelta}
A simple indicator of balance between what the person is asked to do and the support they receive.
\begin{itemize}[leftmargin=*]
  \item \textbf{Purpose.} Keep demands modest and support visible.
  \item \textbf{Typical fields.} Recent asks, recent supports offered, recent supports accepted, user feedback tags, next support suggestion.
  \item \textbf{Produced by.} HEP/ULMS runtime from recent interactions.
  \item \textbf{Consumed by.} Orchestration and equity dashboards.
\end{itemize}

\paragraph{πLīla}
A small, playful practice loop that reinforces skills through quick, varied activities.
\begin{itemize}[leftmargin=*]
  \item \textbf{Purpose.} Make practice sticky without adding burden.
  \item \textbf{Typical fields.} Activity pattern, micro-reward or acknowledgment, reflection prompt, safe repeat limit, cooldown.
  \item \textbf{Produced by.} Content teams and local partners.
  \item \textbf{Consumed by.} HEP/ULMS runtime and S2H micro-brands.
\end{itemize}

\subsection*{De-Identification Governance and Equity Dashboards}
\paragraph{De-ID Governance}
\begin{itemize}[leftmargin=*]
  \item Learning telemetry is de-identified or aggregated by default. Only the minimum needed flows back to drive improvement.
  \item Cohort definitions used for analysis are approved through governance and documented with example size rules to avoid re-identification.
  \item Each export carries a purpose, a retention window, and a receipt that shows fields included and who requested it.
  \item Sensitive free text is avoided. When needed, it is screened for identifiers before storage or export.
\end{itemize}

\paragraph{Equity Dashboards}
\begin{itemize}[leftmargin=*]
  \item Dashboards show access, engagement, comprehension, and time-to-action for approved cohorts.
  \item Gaps are flagged with clear, plain labels. The system proposes concrete adjustments such as alternative media, different timing, or added support.
  \item Trend views compare local results to site-wide baselines so teams see whether changes help close gaps.
  \item Rollout controls let operators slow, pause, or reverse a change until mitigations are in place.
\end{itemize}

\subsection*{Authoring and Release}
\begin{itemize}[leftmargin=*]
  \item Authors work from templates that enforce plain language, accessibility notes, and source links.
  \item Localization captures dialect, cultural references, and reading level adjustments. Changes are versioned and reversible.
  \item A short release note lists what changed, why it matters, and any CSL requirements or guardrails.
\end{itemize}

\subsection*{Interfaces}
\textbf{Inbound.} Triggers and context slices from orchestration, PER preferences, CSLState, and language settings. \\
\textbf{Outbound.} Completion events, comprehension checks, feedback tags, and ReciprocityDelta updates, all de-identified by default.

\subsection*{Glossary}
\begin{itemize}[leftmargin=*]
  \item \textbf{PER:} Personalized Education Record used to tailor delivery and accessibility.
  \item \textbf{MLEM:} Modular Learning Element Model, the atomic lesson or practice unit.
  \item \textbf{TCI:} Task–Content Interface that links a real task to a learning element.
  \item \textbf{CSLState:} Gating state for delivery based on context, safety, and explanation clarity.
  \item \textbf{ReciprocityDelta:} Balance signal between asks and supports to prevent overload.
  \item \textbf{πLīla:} Short, playful practice loop to reinforce skills with variety.
\end{itemize}

\subsection*{Key Performance Indicators (KPIs)}
\begin{longtable}{@{}p{0.26\linewidth}p{0.50\linewidth}p{0.20\linewidth}@{}}
\toprule
\textbf{KPI} & \textbf{Definition} & \textbf{Direction} \\
\midrule
\endhead
Lesson reach & Share of targeted users who receive the lesson within the intended window. & Higher is better \\
Comprehension rate & Share of users who pass a short check or self-report understanding. & Higher is better \\
Action follow-through & Share of users who take the linked task within the expected time. & Higher is better \\
Drop-off rate & Users who start but do not complete a lesson or task. & Lower is better \\
Language match rate & Deliveries in the user’s preferred language or accepted alternative. & Higher is better \\
Accessibility conformance & Lessons that meet declared accessibility needs (captions, contrast, audio). & Higher is better \\
Equity gap index & Largest difference in key outcomes across approved cohorts. & Lower is better \\
Time-to-feedback & Median time from completion to visible update in dashboards. & Lower is better \\
Support uptake & Share of offered supports (reminders, human outreach) that are accepted. & Higher is better \\
Consent clarity score & Users who can state in simple terms what is being collected and why. & Higher is better \\
\bottomrule
\end{longtable}

% End HEP/ULMS Learning Layer

% ---------- Computable Physiology: ADEC–SCM ----------
\section{Computable Physiology: ADEC–SCM}

\subsection*{Role and Mission}
ADEC–SCM is the coordination brain for the loop. It maintains a live picture of a person’s situation, watches for changes, and decides how fast or slow the system should move. Its mission is simple: keep recommendations safe, timely, and proportional while respecting local policy and human oversight.

\subsection*{The State Model in Plain Terms}
Think of the system as tracking a \emph{runtime state} that evolves as new information arrives. This state includes current plans, recent actions, symptoms, access constraints, and support needs. Real life introduces shocks such as a new diagnosis, a missed follow-up, a device failure, or a sudden barrier at home. ADEC–SCM absorbs these shocks, recalculates the situation, and sets a safe pace for next steps. The output is a clear status and a small set of allowed actions with reasons.

\subsection*{Fused Signal and Structural Causal View}
Signals arrive from clinics, apps, devices, and learning activities. ADEC–SCM forms a single \emph{fused signal} that keeps the most useful parts and drops noise. It uses a structural causal view to keep cause and effect straight: which inputs can change outcomes, which are only context, and which should never drive automated action. This approach:
\begin{itemize}[leftmargin=*]
  \item Preserves provenance so teams can see where each piece of evidence came from.
  \item Prefers stable, interpretable features over brittle patterns.
  \item Separates sensitive attributes so they do not steer decisions unless governance allows a narrow, declared use.
  \item Produces explanations that name the few factors that mattered “this time.”
\end{itemize}

\subsection*{Estimation: Keeping a Reliable Picture}
ADEC–SCM keeps its picture of reality fresh using well-known estimation methods. In practical terms:
\begin{itemize}[leftmargin=*]
  \item \textbf{Forecast then correct.} It predicts the near future based on recent behavior, then corrects that prediction when new events arrive.
  \item \textbf{Handle nonlinearity.} When relationships are not straight lines, it still updates sensibly, giving more weight to higher-quality signals.
  \item \textbf{Quantify uncertainty.} Every estimate carries a confidence range so downstream services know when to ask a human to confirm.
  \item \textbf{Fail safe.} If signals conflict or confidence drops, ADEC–SCM slows delivery, asks for a check, or hands off to a person.
\end{itemize}

\subsection*{Control: Choosing the Next Safe Move}
Once the situation is estimated, ADEC–SCM selects the next move using guardrailed optimization:
\begin{itemize}[leftmargin=*]
  \item \textbf{Plan ahead.} It looks a short distance into the future and chooses actions that reduce risk without adding burden.
  \item \textbf{Respect limits.} Clinical rules, local policy, and user preferences act as hard boundaries. The system will not cross them.
  \item \textbf{Learn carefully.} If a learning controller is used, it trains in a simulated space first, starts small, and expands only after review.
  \item \textbf{Explain choices.} Each selected action carries a brief note: why it helps now, what was considered, and when it will be reconsidered.
\end{itemize}

\subsection*{Interfaces and Contracts}
\textbf{Inputs.} Purpose-tagged events from clinical systems and apps, learning completions and feedback, consent states, policy rules, and site settings. \\
\textbf{Outputs.} A small, readable status; allowed action types and pacing; escalation or pause signals; and explanation notes with trace identifiers. \\
\textbf{Contracts.}
\begin{itemize}[leftmargin=*]
  \item Status codes that tools can understand: \textit{normal}, \textit{watch}, \textit{slow}, \textit{pause}, \textit{escalate}.
  \item Pacing guidance such as “one action per day” or “team review before next step.”
  \item Reason fields that list the top drivers and the confidence range.
\end{itemize}

\subsection*{Operating Constraints and Safeguards}
\begin{itemize}[leftmargin=*]
  \item \textbf{Boundary respect.} Protected clinical information stays in clinical space. Learning space receives only what governance allows, typically de-identified aggregates.
  \item \textbf{Fairness checks.} If pacing or escalation differs across groups in ways that raise concern, ADEC–SCM slows or stops automation and opens a review.
  \item \textbf{Kill switch.} Any sign of harm or unstable behavior triggers an immediate stop for the affected feature with a clear recovery path.
  \item \textbf{Version control.} Every change to logic is versioned, reversible, and tied to a release note.
  \item \textbf{Human in the loop.} Sensitive moves require explicit acknowledgment by a clinician or educator before proceeding.
\end{itemize}

\subsection*{Success Criteria}
\begin{itemize}[leftmargin=*]
  \item \textbf{Stability.} Fewer oscillations in recommendations and fewer abrupt escalations without new evidence.
  \item \textbf{Timeliness.} Most safe actions delivered within the expected window after new information arrives.
  \item \textbf{Clarity.} Users report that status and pacing notes are understandable and helpful.
  \item \textbf{Safety.} Near-miss and incident rates trend down; pauses and rollbacks are rare and well-explained.
  \item \textbf{Equity.} Pacing and access to support remain balanced across approved cohorts, or gaps are identified with concrete mitigations.
\end{itemize}

\subsection*{Illustrative Use Cases}
\paragraph{Home blood pressure program}
When morning readings trend late or missing, ADEC–SCM slows new asks, offers a simpler routine, and prompts a human check if patterns persist. If readings stabilize, it restores the usual cadence.

\paragraph{Chemo-related cognitive effects}
When symptom notes suggest worsening fatigue and concentration, ADEC–SCM holds nonessential tasks, surfaces support tips, and opens a team review if red-flag patterns appear.

\paragraph{School-based wellness}
When a class shows low completion during exam week, ADEC–SCM delays delivery, proposes shorter activities, and schedules a catch-up window with teacher approval.

% End Computable Physiology: ADEC–SCM

% ---------- Prime-Indexed Spectral Stack: MQEM/PIRTM ----------
\section{Prime-Indexed Spectral Stack: MQEM/PIRTM}

\subsection*{Role and Scope}
MQEM/PIRTM is the system’s quality and provenance engine. It cleans noisy inputs, guards model and rule stability, detects drift early, and keeps a complete record of what changed, when, and why. It runs quietly in the background yet decides when a release is safe, when to pause, and when to roll back.

\subsection*{Spectral Priors}
Spectral priors are simple expectations about the shape of incoming signals across time and resolution. They help separate meaningful patterns from noise before insights are produced.
\begin{itemize}[leftmargin=*]
  \item \textbf{Structure.} Signals are viewed as a set of frequency-like bands that capture slow trends, daily rhythms, and rapid fluctuations.
  \item \textbf{Use.} Low bands hold long-term habits and seasonality. Mid bands hold weekly cycles and program effects. High bands hold spikes, glitches, and device artifacts.
  \item \textbf{Outcome.} Cleaning happens with restraint so that real changes remain visible and small artifacts fade.
\end{itemize}

\subsection*{PNSM Denoising}
PNSM (Prime-Numbered Spectral Mask) is a conservative denoising procedure that samples and tests non-adjacent bands indexed by prime numbers. It reduces cross-talk between bands and limits overfitting to recent noise.
\begin{itemize}[leftmargin=*]
  \item \textbf{Masking.} Bands that fail basic quality checks are softened or muted, starting with the smallest affected region.
  \item \textbf{Locality.} Adjustments apply only where faults appear, not across the entire signal.
  \item \textbf{Traceability.} Each change records the band tested, the check applied, and the result for later review.
\end{itemize}

\subsection*{Lawful Recursion}
Lawful recursion ensures that any repeated transformation or re-computation follows declared policies and versioned contracts.
\begin{itemize}[leftmargin=*]
  \item \textbf{Bounded steps.} Each pass has a limit on how much it may change a value or threshold.
  \item \textbf{Stop rules.} Re-computation halts when improvement flattens, when confidence drops, or when a safety limit is reached.
  \item \textbf{Human checkpoints.} Sensitive updates require an explicit review note before wider rollout.
\end{itemize}

\subsection*{Drift Probes}
Drift probes are small, fixed datasets and synthetic scenarios used to test for change over time.
\begin{itemize}[leftmargin=*]
  \item \textbf{Coverage.} Probes include typical cases, edge cases, and equity-sensitive cohorts approved by governance.
  \item \textbf{Signals watched.} Data quality, feature stability, threshold movement, false alert rate, and explanation clarity.
  \item \textbf{Actions.} If probes fail, releases pause. The system opens an investigation with suggested checks and a rollback path.
\end{itemize}

\subsection*{Provenance and Lineage}
Provenance answers, in plain terms, how an output came to be.
\begin{itemize}[leftmargin=*]
  \item \textbf{Lineage graph.} Each artifact links to inputs, transforms, owners, dates, and purposes.
  \item \textbf{Receipts.} User-visible receipts summarize the steps that affected their experience and point to sources they are allowed to see.
  \item \textbf{Reproducibility.} Any result can be recreated from pinned versions of data slices, features, rules, and content.
\end{itemize}

\subsection*{Interfaces}
\textbf{Inbound.} Sampled data slices, feature snapshots, candidate rules or content, probe definitions, and policy flags. \\
\textbf{Outbound.} Cleaned data views for testing, drift and quality alerts, release approvals or blocks, lineage records, and user-facing receipts.

\subsection*{Operating Safeguards}
\begin{itemize}[leftmargin=*]
  \item \textbf{Small first.} Every change starts in a shadow space against recent history before any live exposure.
  \item \textbf{Equity lens.} Probe results are broken out by approved cohorts. If gaps widen, the rollout slows and a mitigation plan is required.
  \item \textbf{Rollbacks.} One-click rollback restores the last known good set of features, thresholds, or content with matching receipts.
  \item \textbf{Visibility.} Operators see the exact reason a release is paused, including the probe that failed and the measure that crossed its limit.
\end{itemize}

\subsection*{Initial SLO Targets}
\begin{longtable}{@{}p{0.30\linewidth}p{0.62\linewidth}@{}}
\toprule
\textbf{Area} & \textbf{Target and Notes} \\
\midrule
\endhead
Stability under cleaning & Cleaning changes less than a small fraction of values in any band during normal operation; larger changes trigger review. \\
Drift detection time & Emerging drift is flagged within a short window after threshold crossing on probes or live canaries. \\
Release decision latency & Approve, pause, or roll back within minutes of probe completion for routine changes; human review timelines are posted for sensitive updates. \\
Lineage completeness & Every artifact produced during a change has inputs, owners, dates, and purposes recorded before it becomes eligible for release. \\
Receipt generation rate & User-visible receipts are produced for applicable changes in the vast majority of cases; exceptions are logged with reasons. \\
\bottomrule
\end{longtable}

\subsection*{Success Criteria}
\begin{itemize}[leftmargin=*]
  \item Fewer noisy spikes and fewer false alerts reaching people or teams.
  \item Clear, stable explanations that match what users see.
  \item Fast, precise rollbacks when needed, with minimal user disruption.
  \item Balanced outcomes across approved cohorts while new content or rules are introduced.
\end{itemize}

% End Prime-Indexed Spectral Stack: MQEM/PIRTM

% ---------- One-Loop Reference Dataflow and Contracts ----------
\section{One-Loop Reference Dataflow and Contracts}

\subsection*{Actors and Zones}
This reference view names who sends and receives messages and where protected data is allowed to live. It keeps boundaries visible and makes contracts easy to implement.

\paragraph{Actors}
\begin{itemize}[leftmargin=*]
  \item \textbf{People and Teams} — patients, families, caregivers, clinicians, educators, and operators using apps or portals.
  \item \textbf{S2H Digital Clinic} — turns triggers into checklists, reminders, and team tasks; captures outcomes and barriers.
  \item \textbf{Hologram Health (Record)} — clinical source of truth for summaries, plans, and consents; emits policy-bound events.
  \item \textbf{HEP/ULMS Learning Layer} — delivers short lessons and practice; returns de-identified progress.
  \item \textbf{Intrinsica Data/AI Plane} — ingests events, builds features, registers risk flags, and issues insights and triggers.
  \item \textbf{MQEM/PIRTM Quality and Provenance} — cleans data, detects drift, approves or blocks releases, logs lineage.
  \item \textbf{ADEC--SCM Coordinator} — sets pacing and allowed actions; pauses or escalates when needed.
  \item \textbf{Interop Connectors} — standards-based adapters for clinical and learning systems; webhooks, queues, and filters.
\end{itemize}

\paragraph{Zones}
\begin{itemize}[leftmargin=*]
  \item \textbf{Zone A: Clinical Boundary} — protected health information under clinical governance. Hologram is authoritative here.
  \item \textbf{Zone B: Learning Space} — de-identified or aggregate learning data by default. HEP/ULMS is authoritative here.
  \item \textbf{Zone C: Shared Services} — Intrinsica, MQEM/PIRTM, and ADEC--SCM operating on purpose-tagged events with strict access.
  \item \textbf{Zone D: Edge and Devices} — phones, kiosks, and offline packs with minimal on-device storage and short-lived tokens.
\end{itemize}

\subsection*{Reference Event Sequence (1--9)}
This sequence shows a typical pass through the loop. It is intentionally terse and implementation-agnostic.

\begin{enumerate}[leftmargin=*]
  \item \textbf{Signal captured} (Zone D$\rightarrow$C). A person completes a task, reports a symptom, or a clinic posts a discharge event. Interop attaches purpose and consent references.
  \item \textbf{Normalize and filter} (Zone C). Interop removes unneeded fields, hashes identifiers as required, and stamps a trace identifier.
  \item \textbf{Ingest and feature update} (Zone C). Intrinsica validates the event and updates simple features such as recent activity or follow-up status.
  \item \textbf{Risk and insight} (Zone C). Intrinsica evaluates rules, updates the risk registry, and emits an insight with a plain explanation.
  \item \textbf{Pacing and gating} (Zone C). ADEC--SCM checks status and allows, slows, or pauses next actions based on context and safety.
  \item \textbf{Action orchestration} (Zone C$\rightarrow$D). S2H creates user-facing tasks or team prompts; HEP/ULMS selects a lesson if appropriate.
  \item \textbf{Delivery and acknowledgment} (Zone D). The person sees the action or lesson. The app records completion, barriers, or a help request.
  \item \textbf{Outcome return} (Zone D$\rightarrow$C). Results flow back to Intrinsica. MQEM/PIRTM samples for quality and drift; lineage is updated.
  \item \textbf{Dashboards and review} (Zone C$\rightarrow$A/B). Equity and operations dashboards refresh. Governance may widen, slow, or roll back the change.
\end{enumerate}

\subsection*{Message Contracts}
Contracts are minimal, readable, and versioned. All events share common headers. Bodies vary by topic. Examples below use JSON-like pseudocode for clarity.

\paragraph{Common Headers (all messages)}
\begin{longtable}{@{}p{0.26\linewidth}p{0.64\linewidth}@{}}
\toprule
\textbf{Field} & \textbf{Description} \\
\midrule
\endhead
\texttt{type} & Event name (e.g., \texttt{"risk.update"}, \texttt{"ulms.completion"}). \\
\texttt{version} & Contract version (e.g., \texttt{"1.0"}). \\
\texttt{trace\_id} & End-to-end correlation identifier. \\
\texttt{issued\_at} & ISO~8601 timestamp when the message was created. \\
\texttt{purpose} & Allowed purpose for this processing step (plain terms). \\
\texttt{consent\_ref} & Reference to an active consent or policy basis. \\
\texttt{source} & System of origin (e.g., \texttt{"hologram"}, \texttt{"s2h.app"}). \\
\texttt{site\_id} & Sending site identifier; used for routing and governance. \\
\texttt{idempotency\_key} & Key that makes safe retries possible. \\
\bottomrule
\end{longtable}

\paragraph{\texttt{risk.update} (Intrinsica$\rightarrow$S2H/Record/Governance)}
Announces a change in the risk registry with a human-readable rationale and a safe hint for next actions.

\begin{verbatim}
{
  "type": "risk.update",
  "version": "1.0",
  "trace_id": "trc-7f2b...",
  "issued_at": "2025-10-24T14:03:12Z",
  "purpose": "care-optimization",
  "consent_ref": "cns-48aa...",
  "source": "intrinsica",
  "site_id": "clinic-01",
  "subject_ref": "did:site:xyz:abc123",      // or site pseudonym
  "risk": {
    "id": "bp.adherence",
    "level": "low|moderate|high",
    "status": "new|updated|resolved",
    "reason": "missed 3 of 7 morning readings",
    "expires_at": "2025-11-01T00:00:00Z"
  },
  "explanation": {
    "why_now": "recent misses increased from baseline",
    "visible_to": ["patient","care_team"],
    "learn_more_url": "https://example/why-bp-morning"
  },
  "cs_level": "green|amber|red",
  "next_action_hint": "suggest 2-day routine reset",
  "privacy": { "phi": true, "share_outside_clinical": false }
}
\end{verbatim}

\paragraph{HEP/ULMS Webhooks (HEP/ULMS$\rightarrow$Intrinsica/S2H)}
De-identified by default. Used to keep the loop updated on learning progress without exposing protected clinical details.

\subparagraph{\texttt{ulms.completion}}
\begin{verbatim}
{
  "type": "ulms.completion",
  "version": "1.0",
  "trace_id": "trc-2a1c...",
  "issued_at": "2025-10-24T15:22:44Z",
  "purpose": "learning-outcome",
  "consent_ref": "cns-48aa...",
  "source": "ulms",
  "site_id": "school-09",
  "alias_ref": "pseud-9b7e...",            // no direct PHI; site-limited alias
  "mlem_id": "diet.skin.label-reading.v3",
  "tci_ref": "task-98765",                 // optional: linked real-world task
  "result": {
    "completed": true,
    "duration_sec": 140,
    "comprehension": "pass|fail|skipped"
  },
  "feedback": { "tags": ["too_long","needs_audio"] }
}
\end{verbatim}

\subparagraph{\texttt{ulms.feedback}}
\begin{verbatim}
{
  "type": "ulms.feedback",
  "version": "1.0",
  "trace_id": "trc-2a1c...",
  "issued_at": "2025-10-24T15:23:05Z",
  "purpose": "learning-feedback",
  "consent_ref": "cns-48aa...",
  "source": "ulms",
  "site_id": "school-09",
  "alias_ref": "pseud-9b7e...",
  "mlem_id": "diet.skin.label-reading.v3",
  "rating_1to5": 4,
  "free_text": "examples helped",          // screened to remove identifiers
  "accessibility_used": ["captions"]
}
\end{verbatim}

\subsection*{Transport and Security Expectations}
\begin{itemize}[leftmargin=*]
  \item HTTPS with mTLS or signed tokens. Tokens include site, role, and purpose.
  \item Retries are idempotent. Receivers must ignore exact duplicates.
  \item Webhooks respond \texttt{2xx} on success, \texttt{429} on overload with a \texttt{Retry-After}, and \texttt{4xx/5xx} with reason codes when applicable.
  \item Payloads include only the minimum fields required by the contract. Extra fields are rejected or ignored by policy.
\end{itemize}

\subsection*{Event Acknowledgments}
Two lightweight responses keep back-pressure and observability simple.

\paragraph{\texttt{ack.received}}
\begin{verbatim}
{ "type":"ack.received", "trace_id":"trc-7f2b...", "received_at":"..." }
\end{verbatim}

\paragraph{\texttt{ack.applied}}
\begin{verbatim}
{ "type":"ack.applied", "trace_id":"trc-7f2b...", "applied_at":"...", "artifact_ref":"..." }
\end{verbatim}

\subsection*{Contract Governance}
\begin{itemize}[leftmargin=*]
  \item Versions change only when fields or semantics change. Deprecations carry a clear end date.
  \item Sample payloads accompany every version. Failing examples are added to automated tests.
  \item A public change log covers non-sensitive updates. Detailed lineage stays internal for audits.
\end{itemize}

% End One-Loop Reference Dataflow and Contracts

% ---------- Interoperability and Runtime Safety ----------
\section{Interoperability and Runtime Safety}

\subsection*{Purpose}
Interoperability keeps boundaries clear while allowing systems to cooperate. Runtime safety ensures that recommendations and data flows remain lawful, stable, and understandable during live use.

\subsection*{Standards at a Glance}
\begin{itemize}[leftmargin=*]
  \item \textbf{Clinical: FHIR.} Read and write only what is needed using standard resources (e.g., Patient, CarePlan, Observation, MedicationStatement, DocumentReference, Task). Use purpose tags and audit headers on every call.
  \item \textbf{Learning: LTI and xAPI.} Launch learning items with LTI for secure, role-aware sessions. Capture short outcome statements with xAPI. De-identified by default unless a lawful, approved exception applies.
  \item \textbf{Identity and sessions.} OpenID Connect for sign-in, OAuth~2 for delegated access, short-lived tokens with refresh, and optional mTLS for system-to-system links.
  \item \textbf{Events.} Webhooks for near real-time notifications and a queue for durable delivery with backoff and dead-letter handling.
\end{itemize}

\subsection*{AuthZ and ABAC}
Authorization mixes roles and attributes to keep access narrow and purposeful.
\begin{itemize}[leftmargin=*]
  \item \textbf{Least privilege.} Default to the minimum scope. Expand only with clear purpose and time limits.
  \item \textbf{Role plus attributes.} Roles (clinician, educator, steward, operator) combine with attributes (site, program, cohort, purpose) to decide access at read and write time.
  \item \textbf{Step-up authentication.} Sensitive actions require stronger proof (for example, a second factor or supervised session).
  \item \textbf{Time-boxed admin.} Administrative access expires quickly and is recorded with justification.
  \item \textbf{Scoped tokens.} Tokens carry site, role, and purpose claims; downstream services deny out-of-scope calls.
\end{itemize}

\subsection*{Clinical FHIR Interfaces}
\begin{itemize}[leftmargin=*]
  \item \textbf{Read.} Plans, allergies, medications, key observations, recent documents, and appointments through approved endpoints. Field-level filtering removes unneeded details at the source.
  \item \textbf{Write.} Visit summaries and plan updates as structured documents; Task updates for follow-ups; minimal notes for handoffs.
  \item \textbf{Safeguards.} Per-resource consent checks, rate limits, and idempotency keys. Emergency “break-glass” requires reason codes and auto-review.
\end{itemize}

\subsection*{Learning LTI/xAPI Interfaces}
\begin{itemize}[leftmargin=*]
  \item \textbf{LTI launch.} Role-aware launches with short sessions, redirected back to the Learning Layer when complete.
  \item \textbf{xAPI statements.} Compact completion and comprehension events with alias identifiers and no protected clinical details.
  \item \textbf{Localization.} Language, accessibility, and device capability flags travel with the session so content adapts without extra data sharing.
\end{itemize}

\subsection*{PHI Boundaries}
\begin{itemize}[leftmargin=*]
  \item Clinical records and protected identifiers stay inside clinical systems and Hologram under clinical governance.
  \item Learning systems receive de-identified or aggregate data by default. Any exception is explicit, logged, and time-limited.
  \item Exports include only the minimum fields, with purpose, requester, and retention noted in a receipt.
\end{itemize}

\subsection*{Bedside AI Safety Budgets}
Safety budgets set clear limits on automated activity near the point of care.
\begin{itemize}[leftmargin=*]
  \item \textbf{Frequency limits.} Cap the number of automated prompts or changes per person and per shift.
  \item \textbf{Novelty limits.} Introduce only a small amount of new content at once; require human approval for larger changes.
  \item \textbf{Scope limits.} Reserve sensitive recommendations for human review unless a pathway explicitly allows auto-delivery.
  \item \textbf{Pause rules.} If signals conflict or confidence drops, pause automation and alert a human.
  \item \textbf{Receipts.} Every automated suggestion carries a short “why now” note and a link to sources the user is allowed to see.
\end{itemize}

\subsection*{Auditability}
\begin{itemize}[leftmargin=*]
  \item End-to-end trace identifiers connect input events, decisions, actions, and outcomes.
  \item Append-only audit logs capture who did what, when, for which purpose, and from where.
  \item User-facing receipts summarize data use and are downloadable in readable form.
\end{itemize}

\subsection*{Scheduling and Fairness Guarantees}
Scheduling keeps systems responsive and fair across users and cohorts.
\begin{itemize}[leftmargin=*]
  \item \textbf{No starvation.} Queues guarantee service within posted windows; stale items escalate automatically.
  \item \textbf{Per-user pacing.} Delivery adapts to preference and access constraints so people are not overwhelmed.
  \item \textbf{Cohort fairness.} Dashboards track latency, success, and error rates by approved cohorts. If gaps widen, the system slows expansion and opens a mitigation task.
  \item \textbf{Back-pressure.} When partners are overloaded, the system sheds load gracefully, queues safely, and informs operators.
  \item \textbf{Time-zone and shift awareness.} Delivery windows respect local time, clinic shifts, school hours, and cultural calendars.
\end{itemize}

\subsection*{Runtime Stability and Latency Safeguards}
\begin{itemize}[leftmargin=*]
  \item Health checks and circuit breakers isolate failing components.
  \item Retries use exponential backoff and idempotency keys so duplicates do not occur.
  \item Rate limits and quotas protect partner systems; bursts are buffered.
  \item Degradation paths prefer simple, cached guidance over silence when upstreams are slow.
  \item Rollbacks restore prior versions of rules or content with matching receipts and lineage.
\end{itemize}

\subsection*{Incident Response and Change Control}
\begin{itemize}[leftmargin=*]
  \item Clear severity levels with target response and notice times.
  \item A public change log for non-sensitive updates and an internal log for detailed investigations.
  \item Blameless reviews focused on fixes, documentation, and user communication.
\end{itemize}

% End Interoperability and Runtime Safety

% ---------- Privacy, Consent, and Traceability ----------
\section{Privacy, Consent, and Traceability}

\subsection*{Purpose}
This section defines how the system limits data use to declared purposes, honors consent in practice, and proves what happened after the fact. It sets clear boundaries for what enters each zone, and it provides machine-checkable evidence for oversight, audits, and user rights.

\subsection*{Purpose-Bound Policies}
Policies describe \emph{what data} may be used, \emph{for which purpose}, \emph{by whom}, and \emph{for how long}. Each read or write is evaluated against an active policy.
\begin{itemize}[leftmargin=*]
  \item \textbf{One purpose at a time.} Every operation cites a concrete purpose (for example, \textit{care-transition-checklist}, \textit{patient-education-delivery}, \textit{de-identified-equity-report}).
  \item \textbf{Minimum necessary.} Only the fields required for that purpose are read or written. Extraneous fields are filtered at the boundary.
  \item \textbf{Time scopes.} Policies carry start and end windows. After expiry, the system will not process under that policy without renewal.
  \item \textbf{Role and site scoping.} Access includes role, site, and program constraints. Cross-site access is denied unless a governing policy explicitly allows it.
  \item \textbf{Receipts.} Each approved operation emits a short receipt that references the policy and the purpose used.
\end{itemize}

\subsection*{Consent: Practical, Layered, and Revocable}
Consent is readable, specific, and easy to change.
\begin{itemize}[leftmargin=*]
  \item \textbf{Layered choices.} People see simple on/off switches with links to details. Separate switches exist for learning telemetry, care reminders, research exports, and caregiver sharing.
  \item \textbf{Granular scope.} Consent records include data categories, delivery channels, and allowed partners. Choices are stored with timestamps and the method of capture.
  \item \textbf{Easy revocation.} People can pause or revoke from any device. Revocation stops new processing under that consent and is logged with a receipt.
  \item \textbf{Caregiver and proxy.} Caregivers can be added or removed by the person or by standard clinic processes. All proxy actions generate receipts.
  \item \textbf{Emergency override.} Break-glass use is allowed only under posted conditions, with immediate logging and after-action review.
\end{itemize}

\subsection*{Zero PHI at HEP Ingress}
The education space (HEP/ULMS) runs on de-identified or aggregate data by default.
\begin{itemize}[leftmargin=*]
  \item \textbf{Boundary rule.} Protected clinical identifiers do not enter HEP/ULMS. Sessions use site-limited aliases and short-lived tokens.
  \item \textbf{Field filtering.} Interop removes direct identifiers and drops unnecessary context before delivery to HEP/ULMS.
  \item \textbf{Cohort safeguards.} Cohorts used for analysis follow minimum-size rules to reduce re-identification risk.
  \item \textbf{Exception handling.} Rare, lawful exceptions require explicit policy, a time limit, and a dedicated receipt trail. Defaults return to de-identified flow when the exception expires.
\end{itemize}

\subsection*{Ledgered Provenance}
Every important step produces an immutable, append-only record.
\begin{itemize}[leftmargin=*]
  \item \textbf{What is recorded.} Inputs referenced, purpose used, decision made, user or service actor, timestamps, and any artifacts produced (for example, a plan update or a lesson assignment).
  \item \textbf{How it is used.} Operators reconstruct timelines, auditors verify constraints, and users can download plain-language summaries of their own activity.
  \item \textbf{Versioning.} Content, rules, and policies are versioned. Releases are signed, and rollbacks point to prior signed versions.
\end{itemize}

\subsection*{Machine-Checkable Audits}
Audits rely on structured evidence that tools can verify without guesswork.
\begin{itemize}[leftmargin=*]
  \item \textbf{Fixed schemas.} Audit entries follow fixed field names and types so queries remain stable over time.
  \item \textbf{Linked traces.} A single trace identifier connects intake, decisions, actions, and outcomes.
  \item \textbf{Programmatic tests.} Routine jobs check for missing receipts, expired policies in use, or unauthorized fields crossing a boundary.
  \item \textbf{Export packs.} Sites can generate time-bounded audit packs that include logs, policy snapshots, and example receipts for regulators or IRBs.
\end{itemize}

\subsection*{Receipts People Can Read}
Receipts translate system events into plain language and are accessible from apps and portals.
\begin{itemize}[leftmargin=*]
  \item \textbf{Core elements.} What happened, why it happened, which data moved, who could see it, and how to change future use.
  \item \textbf{Delivery.} Receipts appear in-app and are available as downloadable files. Paper copies can be printed on request at clinics or schools.
  \item \textbf{Language access.} Receipts are localized and meet accessibility guidelines (readable fonts, captions for media, and screen-reader friendly structure).
\end{itemize}

\subsection*{Reference States and Records}
\paragraph{Consent states}
\begin{longtable}{@{}p{0.22\linewidth}p{0.70\linewidth}@{}}
\toprule
\textbf{State} & \textbf{Meaning} \\
\midrule
\endhead
\textit{active} & Processing allowed within the stated purpose, time window, and channels. \\
\textit{paused} & Temporarily suspended by the person or site; no new processing under this consent. \\
\textit{revoked} & Ended by the person; future processing blocked. Historical uses remain logged. \\
\textit{expired} & Time window ended; renewal required. \\
\textit{emergency\_override} & Limited, documented use under posted emergency policy with after-action review. \\
\bottomrule
\end{longtable}

\paragraph{Audit record fields}
\begin{longtable}{@{}p{0.28\linewidth}p{0.64\linewidth}@{}}
\toprule
\textbf{Field} & \textbf{Description} \\
\midrule
\endhead
\texttt{trace\_id} & End-to-end correlation ID. \\
\texttt{actor} & User or service that performed the action, with role and site. \\
\texttt{purpose} & Declared purpose that allowed the action. \\
\texttt{policy\_ref} & Versioned policy identifier used to approve the action. \\
\texttt{data\_scope} & Categories and fields accessed or written. \\
\texttt{timestamp} & ISO~8601 time when the event occurred. \\
\texttt{outcome} & Accepted, denied, paused, or escalated with reason. \\
\texttt{receipt\_ref} & Pointer to the human-readable receipt when applicable. \\
\bottomrule
\end{longtable}

\subsection*{User Rights and Self-Service}
\begin{itemize}[leftmargin=*]
  \item \textbf{See and change.} People can view consents, change choices, and download receipts from any supported device.
  \item \textbf{Access and export.} People can request a copy of their information in a portable format with clear explanations.
  \item \textbf{Rectify.} People can flag inaccuracies in summaries or plans for review and correction.
  \item \textbf{Complain.} People can submit concerns. The system logs the case, response time, and outcome.
\end{itemize}

\subsection*{Operating Commitments}
\begin{itemize}[leftmargin=*]
  \item \textbf{No silent expansion.} New data uses require an updated policy and change notice; sensitive changes require explicit re-consent.
  \item \textbf{Independent review.} Sites maintain a standing privacy and ethics review function that samples receipts and audits regularly.
  \item \textbf{Graceful degradation.} If a boundary fails closed (for example, token problems), the system favors withholding data over risking misuse and informs operators.
\end{itemize}

\subsection*{Success Criteria}
\begin{itemize}[leftmargin=*]
  \item High share of operations with valid purpose and consent references.
  \item Low rate of boundary violations or missing receipts.
  \item Short resolution time for consent changes and data access requests.
  \item Positive user ratings on clarity of notices and receipts.
\end{itemize}

% End Privacy, Consent, and Traceability


% ---------- SLOs, KPIs, and Equity ----------
\section{SLOs, KPIs, and Equity}

\subsection*{Purpose}
This section defines how the system proves it is dependable in daily use and how sites know whether it is helping people fairly. Service Level Objectives (SLOs) describe the expected behavior of the live system. Key Performance Indicators (KPIs) describe changes in human outcomes and team workflow. Equity measures ensure that benefits reach all approved cohorts and that any gaps are seen early and addressed.

\subsection*{Runtime SLO Invariants}
These invariants apply to every site. They are built into the platform and do not depend on local configuration.
\begin{itemize}[leftmargin=*]
  \item \textbf{Traceability always on.} Every message and action carries a trace identifier and a readable reason.
  \item \textbf{Minimum necessary data.} Events and payloads include only fields allowed for the stated purpose.
  \item \textbf{Safe retries.} Idempotency keys prevent duplicates during retry and recovery.
  \item \textbf{Graceful degradation.} If an upstream is slow, the system serves cached guidance or pauses rather than sending incomplete or risky suggestions.
  \item \textbf{Human override.} Any user with the right role can pause automation for a person or a cohort with a single control, generating a receipt.
  \item \textbf{Kill switch.} Operators can disable a feature or release immediately when harm or instability appears.
  \item \textbf{Equity guard.} When equity probes fail, expansion stops automatically until a mitigation is in place.
\end{itemize}

\subsection*{Site SLO Targets}
Targets are tuned with each site and posted on an operations dashboard. Examples below show a starting set for first deployments.

\begin{longtable}{@{}p{0.28\linewidth}p{0.56\linewidth}p{0.12\linewidth}@{}}
\toprule
\textbf{Area} & \textbf{Target at Launch} & \textbf{Window} \\
\midrule
\endhead
Action delivery latency & Most user-facing actions appear within seconds after a triggering event is received and validated. & P95 \\
Data freshness & New clinical and learning events are visible to downstream services within minutes, subject to partner system cadence. & P95 \\
Availability & Core runtime services available during posted service hours, with planned maintenance windows announced in advance. & Monthly \\
Audit completeness & Nearly all events and decisions include trace ID, purpose, and consent reference; missing entries trigger alerts. & \(\geq\) 99\% \\
Explanation coverage & Insights and recommendations include a plain “why this, why now” note. & \(\geq\) 98\% \\
Issue response & Confirmed safety or privacy issues triaged immediately with user notice following site policy. & On receipt \\
Privacy pass-rate (runtime) & Processing attempts that pass policy and consent checks without manual intervention. & \(\geq\) 99\% \\
\bottomrule
\end{longtable}

\subsection*{KPI Definitions}
KPIs are plain, stable measures with clear methods. They focus on whether people understand and act, whether risk changes in useful ways, and whether teams acknowledge and use the system.

\paragraph{Engagement}
Share of targeted users who open or view at least one delivered action or lesson within the intended window. Count one engagement per window per user to avoid inflation.

\paragraph{Adherence}
Share of delivered tasks that users complete within the expected time. Report separately for simple self-care steps, learning items, and team tasks that require a human response.

\paragraph{Risk-Delta}
Change in active risk flags over a site-defined period. Count resolved or downgraded items as positive movement and newly raised or upgraded items as negative movement. Present both direction and magnitude so teams see whether change is meaningful.

\paragraph{Equity Gap}
Largest absolute difference in a KPI (for example, adherence) between any two approved cohorts during the same period. Cohorts are defined in governance with minimum-size rules. Gaps trigger a review when they exceed posted thresholds.

\paragraph{Privacy Pass-Rate}
Share of processing attempts that proceed under a valid purpose and consent reference without manual exception. Report by flow (clinical, learning, research export) and by site area. Investigate any drops quickly.

\paragraph{Provider Acknowledgment}
Share of insight cards or team tasks that receive a human acknowledgment within the expected window. Track both first acknowledgment and time to close. Use this to tune volume, timing, and clarity.

\subsection*{Dashboards and Breakouts}
Dashboards present live SLOs and KPIs together so teams see cause and effect.
\begin{itemize}[leftmargin=*]
  \item \textbf{Live status.} Availability, delivery latency, data freshness, backlog depth, and error classes with recent incidents and mitigations.
  \item \textbf{Outcome tiles.} Engagement, adherence, risk-delta, provider acknowledgment, and privacy pass-rate.
  \item \textbf{Equity view.} Side-by-side charts by approved cohorts with clear labels and thresholds. Gaps show a red or amber badge until closed.
  \item \textbf{Drill-down.} Click any tile to see definitions, recent changes, and the content or pathways most associated with the result.
\end{itemize}

\subsection*{Target-Setting Method}
Targets start modest, then tighten with evidence.
\begin{itemize}[leftmargin=*]
  \item \textbf{Baseline and ramp.} Measure two to four weeks of baseline, then set a first ramp that is achievable without strain.
  \item \textbf{Cohort-aware.} Post targets separately for cohorts when appropriate so improvements are visible across groups, not only on averages.
  \item \textbf{Content-linked.} Tie changes in KPIs to content or pathway versions to see what works and what should pause.
  \item \textbf{Public notes.} Post short notes for each target explaining why it matters and how it will be reviewed.
\end{itemize}

\subsection*{Equity Actions}
Equity is an operating rule, not an afterthought.
\begin{itemize}[leftmargin=*]
  \item \textbf{Acceptance bands.} Define green, amber, and red bands for equity gaps. Amber opens a mitigation task; red pauses expansion or triggers rollback.
  \item \textbf{Mitigation kits.} Offer standard responses such as alternate media, different delivery windows, human outreach, or community partner support.
  \item \textbf{Local context.} Invite community input when gaps persist. Adjust content and channels to match local realities.
  \item \textbf{Re-test.} After mitigation, require a re-measure before resuming expansion.
\end{itemize}

\subsection*{Data Quality and Missingness}
\begin{itemize}[leftmargin=*]
  \item \textbf{Completeness flags.} KPIs display a data-quality badge so readers know when results rely on small samples or missing fields.
  \item \textbf{Conservative defaults.} When key data are missing, dashboards show “insufficient data” rather than a misleading number.
  \item \textbf{Attribution notes.} Results identify which content or pathway versions were in play during the period.
\end{itemize}

\subsection*{Review Cadence and Governance}
\begin{itemize}[leftmargin=*]
  \item \textbf{Weekly ops review.} Check SLOs, backlog, incidents, and privacy pass-rate; assign fixes with owners and dates.
  \item \textbf{Biweekly equity review.} Inspect cohort gaps, mitigation status, and any paused expansions.
  \item \textbf{Monthly outcomes review.} Examine trends in engagement, adherence, risk-delta, and provider acknowledgment; update targets as needed.
  \item \textbf{Quarterly public summary.} Publish non-sensitive results and explain what changed, why, and what is next.
\end{itemize}

\subsection*{Success Criteria}
\begin{itemize}[leftmargin=*]
  \item Runtime SLOs hold steady during peak periods with quick recovery from faults.
  \item Engagement and adherence improve without increasing alert fatigue.
  \item Risk-delta trends positive for targeted pathways.
  \item Provider acknowledgment stays high with acceptable time to close.
  \item Equity gaps narrow or remain small across approved cohorts.
  \item Privacy pass-rate remains high with few exceptions and fast resolution.
\end{itemize}

% End SLOs, KPIs, and Equity

% ---------- Example Care–Learning Pathways ----------
\section{Example Care–Learning Pathways}

\subsection*{ICU “Sepsis Watch” Loop and Contracts}
\paragraph{Purpose and scope}
Detect suspected sepsis early in the ICU, focus attention, and make follow-through simple. Keep alerts few, clear, and accountable. Pair every alert with a short checklist and a learning bite that fits the moment.

\paragraph{Trigger sources}
Lab deltas, vital sign changes, antimicrobial orders, nurse concern notes, and device signals that meet site-approved criteria for “suspected sepsis.” Triggers carry purpose tags and a trace identifier.

\paragraph{Roles at a glance}
\begin{longtable}{@{}p{0.22\linewidth}p{0.74\linewidth}@{}}
\toprule
\textbf{Role} & \textbf{Responsibilities in this loop} \\
\midrule
\endhead
Bedside nurse & Receives alert, runs the short checklist, documents findings, flags barriers, and escalates if needed. \\
ICU physician / APP & Reviews summary, confirms or rules out sepsis, orders next steps, and acknowledges final status. \\
Pharmacist & Verifies antimicrobial choice and timing when the checklist calls for it. \\
Charge nurse / coordinator & Monitors loop health, reassigns tasks if delays occur. \\
\bottomrule
\end{longtable}

\paragraph{Event sequence (high level)}
\begin{enumerate}[leftmargin=*]
  \item \textbf{Suspect signal detected.} Intrinsica registers a sepsis-suspect risk with a plain explanation and confidence band.
  \item \textbf{Pacing check.} ADEC–SCM confirms it is safe to alert now; if not, it slows or routes to team review.
  \item \textbf{Action creation.} S2H creates a one-page \emph{Sepsis Watch} checklist and a team acknowledgment task.
  \item \textbf{Delivery.} The nurse console and on-call device receive the task; audible or silent mode follows unit policy.
  \item \textbf{At-the-bedside steps.} Nurse completes quick checks (fluids status, lactate draw status, cultures if ordered, antimicrobial timing) and notes barriers.
  \item \textbf{Learning bite.} ULMS offers a 60–90 second refresher on timing bundles, with captions and quick graphics.
  \item \textbf{Team acknowledgment.} Physician/APP reviews the auto-generated summary card and taps \emph{confirm}, \emph{defer}, or \emph{rule out}, with a short reason.
  \item \textbf{Outcome return.} Results flow back to Intrinsica; MQEM/PIRTM samples for quality and explanation clarity.
  \item \textbf{Close or escalate.} If items remain open past the window, S2H escalates to charge nurse; otherwise the loop closes with a receipt.
\end{enumerate}

\paragraph{Minimal message contracts (pseudocode)}
Only message bodies are shown; common headers (type, version, trace\_id, purpose, consent\_ref, site\_id) apply.

\subparagraph{\texttt{risk.update} (sepsis suspect)}
\begin{verbatim}
{
  "risk": {
    "id": "sepsis.suspect",
    "level": "moderate",
    "status": "new",
    "reason": "lactate rise + hypotension + antibiotic order",
    "expires_at": "2025-10-24T18:00:00Z"
  },
  "explanation": {
    "why_now": "pattern crossed site threshold",
    "visible_to": ["care_team"]
  },
  "cs_level": "amber",
  "next_action_hint": "start Sepsis Watch checklist"
}
\end{verbatim}

\subparagraph{\texttt{s2h.task.create} (Sepsis Watch checklist)}
\begin{verbatim}
{
  "task": {
    "id": "task-sepsiswatch-4129",
    "title": "Sepsis Watch - 6 items",
    "due_in_min": 30,
    "assignee_role": "icu_nurse",
    "checklist": [
      {"id":"fluids.status","label":"Fluids status confirmed"},
      {"id":"lactate.draw","label":"Lactate drawn/repeat ordered"},
      {"id":"cultures","label":"Cultures queued if ordered"},
      {"id":"abx.timer","label":"Antibiotic timing within goal"},
      {"id":"organ.notes","label":"Organ function notes"},
      {"id":"notify.app","label":"APP/MD notified"}
    ]
  }
}
\end{verbatim}

\subparagraph{\texttt{ulms.assignment} (bundle refresher)}
\begin{verbatim}
{
  "lesson": {
    "mlem_id": "sepsis.bundle.timeliness.v2",
    "est_time_sec": 75,
    "accessibility": ["captions","high_contrast"]
  }
}
\end{verbatim}

\subparagraph{\texttt{s2h.task.outcome}}
\begin{verbatim}
{
  "task_id": "task-sepsiswatch-4129",
  "completed": true,
  "barriers": ["lab_delay"],
  "notes": "antibiotic hung at 16:42, repeat lactate queued"
}
\end{verbatim}

\subparagraph{\texttt{record.summary.add}}
\begin{verbatim}
{
  "summary": {
    "title": "Sepsis Watch outcome",
    "body": "Checklist complete; MD confirmed 'suspected'. Next: repeat lactate."
  }
}
\end{verbatim}

\paragraph{Operating constraints and success checks}
\begin{itemize}[leftmargin=*]
  \item One alert per patient within a cooldown window unless status worsens materially.
  \item Target nurse start time within minutes of alert; target physician acknowledgment within posted ICU window.
  \item Learning bite is optional when workload is high; the system defers it automatically if tasks run long.
  \item Equity view includes language access and shift patterns to detect gaps in completion and acknowledgment.
\end{itemize}

\subsection*{Heart Failure (HF) Transitions Loop and Contracts}
\paragraph{Purpose and scope}
Reduce avoidable readmissions by making discharge and the first weeks at home clear and supported. Focus on medications, weight monitoring, sodium guidance, and early follow-up.

\paragraph{Milestones and actions}
\begin{longtable}{@{}p{0.20\linewidth}p{0.76\linewidth}@{}}
\toprule
\textbf{Milestone} & \textbf{Core actions and artifacts} \\
\midrule
\endhead
Pre-discharge & Confirm meds list and doses; schedule follow-up; provide plain-language plan with warning signs; add caregiver contact if consented. \\
Day 2–3 check-in & Text or call prompt; confirm meds in hand; weight monitoring setup; barriers capture (cost, device, transport). \\
Day 7 visit & Team review of weights and symptoms; adjust plan; reinforce sodium and fluid guidance. \\
Day 30 wrap-up & Evaluate adherence, events, and learning completion; update long-term plan. \\
\bottomrule
\end{longtable}

\paragraph{Event sequence (high level)}
\begin{enumerate}[leftmargin=*]
  \item \textbf{Discharge ready.} Hologram posts a concise discharge summary with meds and follow-up.
  \item \textbf{Starter pack.} S2H assembles an \emph{HF Start} pack: weight card, sodium guide, meds checklist, and a caregiver share link if allowed.
  \item \textbf{At-home rhythm.} Daily weight reminder in the morning; symptom check twice a week; quick help route always visible.
  \item \textbf{Learning moments.} ULMS assigns two-minute lessons: reading labels, scale placement, and when to call.
  \item \textbf{Barrier handling.} If weights are missing or meds not in hand, S2H opens a social support or pharmacy task.
  \item \textbf{Clinic sync.} Before the Day 7 visit, S2H compiles a one-page summary for the team.
  \item \textbf{Outcome review.} Intrinsica updates risk; MQEM/PIRTM checks for false alerts and messaging volume.
\end{enumerate}

\paragraph{Minimal message contracts (pseudocode)}
\subparagraph{\texttt{s2h.task.create} (HF starter pack)}
\begin{verbatim}
{
  "task": {
    "id": "task-hfstart-1001",
    "title": "HF Start - home setup",
    "assignee_role": "patient",
    "items": [
      {"id":"scale.setup","label":"Place scale and test reading"},
      {"id":"meds.inhand","label":"Confirm meds in hand"},
      {"id":"followup.verify","label":"Verify clinic follow-up time"}
    ],
    "due_in_hours": 24
  }
}
\end{verbatim}

\subparagraph{\texttt{ulms.assignment} (two-minute lessons)}
\begin{verbatim}
{
  "lesson": {
    "mlem_id": "hf.labels.sodium.v1",
    "est_time_sec": 120,
    "language": "preferred"
  }
}
\end{verbatim}

\subparagraph{\texttt{s2h.barrier.report}}
\begin{verbatim}
{
  "task_id": "task-hfstart-1001",
  "barrier": "no_scale",
  "help_request": true
}
\end{verbatim}

\subparagraph{\texttt{risk.update} (adherence trend)}
\begin{verbatim}
{
  "risk": {
    "id": "hf.monitoring.gap",
    "level": "low",
    "status": "updated",
    "reason": "missing 2 of last 3 weights"
  },
  "next_action_hint": "offer call from coach"
}
\end{verbatim}

\paragraph{Clinic-ready summaries}
\begin{itemize}[leftmargin=*]
  \item One-page \emph{HF Week 1} card: weight sparkline, missed days, self-reported symptoms, meds barriers, learning completion.
  \item Plain recommendations: “keep,” “adjust,” or “review,” each with a short reason and next step.
\end{itemize}

\subsection*{Education Actions and SLO Gates Across Pathways}
\paragraph{Education actions}
\begin{itemize}[leftmargin=*]
  \item \textbf{Sepsis Watch.} Optional 60–90 second refresher tied to timing goals; skippable when workload is high.
  \item \textbf{HF Transitions.} Two-minute lessons on sodium, daily weights, red flags, and med lists; audio and captions default on.
  \item \textbf{Barrier-aware nudges.} When people skip actions, the system asks why and offers an alternative or opens a help route.
\end{itemize}

\paragraph{SLO gates}
\begin{itemize}[leftmargin=*]
  \item \textbf{Delivery.} ICU alerts and bedside tasks appear within the posted window; HF morning reminders arrive in the user’s local morning.
  \item \textbf{Acknowledgment.} Sepsis Watch requires a clinician acknowledgment within the ICU target window; HF Day 2–3 check-ins require staff review if barriers appear.
  \item \textbf{Privacy pass-rate.} Learning assignments in both pathways use de-identified telemetry; exceptions are visible on the dashboard with expiry dates.
  \item \textbf{Equity.} Dashboards break out completion and acknowledgment by approved cohorts. Red-band gaps pause expansion until mitigations land.
\end{itemize}

\paragraph{Success criteria}
\begin{itemize}[leftmargin=*]
  \item \textbf{Sepsis Watch.} High acknowledgment within window, fewer false alerts, faster bundle timing, clear documentation.
  \item \textbf{HF Transitions.} Higher day-7 visit readiness, steady weight monitoring, fewer unplanned returns, clear user feedback on usefulness.
  \item \textbf{Learning.} Short lessons completed at high rates; people can explain “what to do next” in their own words.
\end{itemize}

% End Example Care–Learning Pathways


% ---------- Neurodiversity and Inclusion ----------
\section{Neurodiversity and Inclusion}

\subsection*{Purpose}
This section ensures the system is welcoming and usable for neurodivergent people, including those with autism, ADHD, dyslexia, dyspraxia, sensory processing differences, and related profiles. It defines \emph{EchoBraid} controls for stimulus, language, and pacing, sets concrete acceptance criteria, and explains equity acceptance bands and privacy gates specific to this space.

\subsection*{EchoBraid: Controls People Can Set}
EchoBraid is a compact control panel that travels with the user across apps and moments. It lets people shape how the system “speaks,” “shows,” and “waits,” without extra forms.

\paragraph{Stimulus controls}
\begin{itemize}[leftmargin=*]
  \item \textbf{Visual load:} simple / medium / rich. Controls animation, clutter, and contrast.
  \item \textbf{Motion:} reduce motion on / off. Disables non-essential movement.
  \item \textbf{Audio:} captions default on / off; voice hints on / off; volume range.
  \item \textbf{Color sensitivity:} avoid flashing and high-saturation palettes on / off.
\end{itemize}

\paragraph{Language controls}
\begin{itemize}[leftmargin=*]
  \item \textbf{Reading mode:} short bullets / short narrative / audio first.
  \item \textbf{Vocabulary:} plain terms only / plain + definitions on hover.
  \item \textbf{Structure:} numbered steps before background / background first.
\end{itemize}

\paragraph{Pacing and time}
\begin{itemize}[leftmargin=*]
  \item \textbf{Task size:} micro-steps (≤2 minutes) / standard steps (≤5 minutes).
  \item \textbf{Timer:} visible progress bar on / off; extend time on request.
  \item \textbf{Breaks:} built-in pause prompts after N steps.
\end{itemize}

\paragraph{Executive function supports}
\begin{itemize}[leftmargin=*]
  \item \textbf{Checklist mode:} always show checkbox lists and one clear next step.
  \item \textbf{Preview:} show “what will happen” before starting any flow.
  \item \textbf{Saved state:} resume exactly where the person left off, without penalty.
\end{itemize}

\paragraph{Social-context and interaction}
\begin{itemize}[leftmargin=*]
  \item \textbf{Directness:} direct requests / suggestions only (no figurative language).
  \item \textbf{Masking-aware tips:} optional scripts for calls, visits, or forms.
  \item \textbf{Noise-safe mode:} enlarge touch targets and simplify screens for busy settings.
\end{itemize}

\paragraph{Notification style}
\begin{itemize}[leftmargin=*]
  \item \textbf{Frequency:} quiet / standard / coaching. Caps repeat nudges per day.
  \item \textbf{Channel:} SMS / in-app / email, with quiet hours by local time.
  \item \textbf{Acknowledge without typing:} one-tap yes/no and quick-reasons.
\end{itemize}

\paragraph{Data display}
\begin{itemize}[leftmargin=*]
  \item \textbf{Cognitive load badges:} show an estimate of reading time and steps.
  \item \textbf{Consistency:} same place for back, next, save, and help across apps.
  \item \textbf{Sensory warnings:} brief notice before any content with intense visuals or audio.
\end{itemize}

\subsection*{Acceptance Criteria for EchoBraid}
These are testable checks used during pilots and releases.

\paragraph{Usability}
\begin{itemize}[leftmargin=*]
  \item A person can find EchoBraid within two taps or clicks from any screen.
  \item Changing any control updates the current screen within two seconds.
  \item “Restore defaults” returns to site defaults without losing saved progress.
  \item All controls are keyboard- and screen-reader-accessible.
\end{itemize}

\paragraph{Clarity}
\begin{itemize}[leftmargin=*]
  \item Each control has a one-line label and a short “learn more” that fits on one small card.
  \item Cognitive load badges are visible before a task starts and match actual step counts within a small margin.
\end{itemize}

\paragraph{Pacing}
\begin{itemize}[leftmargin=*]
  \item Micro-step mode ensures no single step exceeds the posted time estimate.
  \item Break prompts appear at the configured interval and never interrupt a confirmation tap.
\end{itemize}

\paragraph{Consistency}
\begin{itemize}[leftmargin=*]
  \item The location and wording of \emph{Back}, \emph{Next}, \emph{Save}, and \emph{Help} do not change across micro-brands.
  \item Notification quiet hours are respected with no more than one exception per month for safety-critical items.
\end{itemize}

\paragraph{Error recovery}
\begin{itemize}[leftmargin=*]
  \item If the app reloads, the person returns to the same step with no data loss.
  \item If a device goes offline, progress can be saved locally and synced later with a clear indicator.
\end{itemize}

\subsection*{Equity Acceptance Bands}
Equity bands help teams decide when to go faster, hold steady, or slow down for specific cohorts in neurodiversity-related flows. Bands are computed over a rolling window and shown on dashboards.

\begin{longtable}{@{}p{0.18\linewidth}p{0.50\linewidth}p{0.26\linewidth}@{}}
\toprule
\textbf{Band} & \textbf{Meaning in practice} & \textbf{Required action} \\
\midrule
\endhead
\textbf{Green} & Access, completion, and comprehension within small variance of site baseline. & Continue rollout; monitor weekly. \\
\textbf{Amber} & Noticeable gap in one or two measures (e.g., completion lower, time-to-complete longer). & Add supports (audio, micro-steps, human outreach); reassess in one week. \\
\textbf{Red} & Sustained gap across multiple measures or safety complaints. & Pause expansion for affected cohorts; run fixes; resume only after re-test meets Green/Amber. \\
\bottomrule
\end{longtable}

\paragraph{Neurodiversity-focused measures}
\begin{itemize}[leftmargin=*]
  \item \textbf{Sensory opt-in rate:} share of users setting at least one stimulus control.
  \item \textbf{Micro-step preference rate:} share selecting micro-steps and their completion rate.
  \item \textbf{Quiet-hour delivery success:} deliveries outside quiet hours as a share of total.
  \item \textbf{Time-to-first-success:} median time from first view to first completed step.
  \item \textbf{Help route uptake:} share of users who accept a human assist when offered.
\end{itemize}

\subsection*{Privacy Gates for Neurodiversity Features}
Privacy gates are explicit checks that run before sensitive features are used.

\paragraph{Gate categories}
\begin{itemize}[leftmargin=*]
  \item \textbf{Preference-only gate:} EchoBraid settings are treated as preferences, not diagnoses. Stored with minimal fields and short retention.
  \item \textbf{Sensitive media gate:} any request to capture or process voice, image, or video requires a separate, time-limited consent specific to that use.
  \item \textbf{Context gate:} school or home context cannot be inferred or shared unless a policy explicitly allows it with a posted purpose.
\end{itemize}

\paragraph{What the system must do}
\begin{itemize}[leftmargin=*]
  \item Show a plain explanation before activating any sensitive feature; proceed only after an explicit yes.
  \item Offer a non-sensitive alternative (e.g., text instead of voice) without penalty.
  \item Generate a short, readable receipt that lists what was used, for what purpose, and how to change it later.
  \item Honor revocation immediately and revert to default, non-sensitive modes.
\end{itemize}

\subsection*{Interfaces and Handshakes}
\paragraph{Inbound (from user and team tools)}
EchoBraid preference updates; help requests; quiet-hour schedules; language and accessibility hints.

\paragraph{Outbound (to delivery and dashboards)}
Adjusted content variants (micro-steps, captions, low-motion), success markers, time-to-complete, help uptake, and equity band status for approved cohorts.

\subsection*{Success Criteria}
\begin{itemize}[leftmargin=*]
  \item People report that screens are predictable, calm, and easy to follow.
  \item Completion and comprehension improve for users who enable EchoBraid compared to their own baseline.
  \item Equity bands move from Red or Amber toward Green after supports are added.
  \item Sensitive features are used only behind explicit gates with clear receipts.
\end{itemize}

% End Neurodiversity and Inclusion
% ---------- Failure Modes and Compensations ----------
\section{Failure Modes and Compensations}

\subsection*{Purpose}
This section defines how the system behaves when things go wrong. It lists common failure modes, how they are detected, the automatic compensations that run immediately, and the follow-up actions for people and teams. The goal is simple: lose nothing important, avoid duplicate work, and recover quickly with a clear record of what happened.

\subsection*{Failure Mode Catalog}
\begin{longtable}{@{}p{0.22\linewidth}p{0.28\linewidth}p{0.26\linewidth}p{0.20\linewidth}@{}}
\toprule
\textbf{Failure mode} & \textbf{Detection} & \textbf{Automatic compensation} & \textbf{Human follow-up} \\
\midrule
\endhead
Webhook endpoint unavailable & HTTP timeouts or 5xx & Retry with backoff; write to durable queue; raise low-severity alert & Ops checks partner status page; coordinate reopen time \\
Duplicate webhook delivery & Repeated \texttt{idempotency\_key} & De-duplicate; return \texttt{200 OK} without reprocessing & None \\
Message order confusion & Gaps in \texttt{trace\_id} sequence & Reorder using timestamps; hold later messages briefly & Review stuck items; clear backlog rule if needed \\
Model or rule regression & Probe failure or spike in errors & Pause release; route to prior signed version & Owner posts fix note; schedule re-test \\
Corrupt or missing content asset & 4xx/5xx on asset fetch & Swap to cached safe copy; mark item as degraded & Content team restores asset; rerun checks \\
Offline edits collide & Conflicting field edits on sync & CRDT merge where safe; open human reconciliation view & Staff confirms final value; receipt issued \\
Clock or timezone drift & Implausible timestamps & Normalize to server time; flag for review & Site adjusts device time settings \\
\bottomrule
\end{longtable}

\subsection*{Webhook Retries}
\paragraph{Principles}
Never drop an event. Never process it twice. Keep partners safe when they are overloaded.

\paragraph{Behavior}
\begin{itemize}[leftmargin=*]
  \item \textbf{Backoff tiers.} Quick retries for transient faults, then slower intervals, then a dead-letter queue (DLQ) if exhaustion occurs. Operators can replay from DLQ with one click after a fix.
  \item \textbf{Idempotency.} Every message carries an \texttt{idempotency\_key}. Receivers must store recent keys and return success on duplicate deliveries without side effects.
  \item \textbf{Signals in headers.} Senders include reason codes for retries and the next attempt time; receivers respond with clear \texttt{2xx/4xx/5xx/429} and, when relevant, \texttt{Retry-After}.
  \item \textbf{Back-pressure.} On \texttt{429}, the sender slows the stream and widens retry windows to protect the receiver.
  \item \textbf{Partial outage handling.} If a downstream is down, events for that partner are queued while other partners continue normally.
\end{itemize}

\paragraph{Example retry envelope (pseudocode)}
\begin{verbatim}
Headers:
  Idempotency-Key: 08b1-...-9f
  Trace-Id: trc-55a2...
  Purpose: care-optimization
  Retry-Attempt: 3
  Next-Retry-At: 2025-10-24T16:05:00Z
Body:
  { "type": "risk.update", "version": "1.0", ... }
\end{verbatim}

\subsection*{Model and Rule Rollback}
\paragraph{Principles}
Changes must be reversible within minutes. People should not feel turbulence when a release misbehaves.

\paragraph{Behavior}
\begin{itemize}[leftmargin=*]
  \item \textbf{Signed releases.} Content, thresholds, features, and rules ship as versioned bundles with signatures and short release notes.
  \item \textbf{Canaries first.} New bundles run on a small slice or in shadow mode. MQEM/PIRTM probes watch stability, volume, and equity impact.
  \item \textbf{Instant pause.} Any probe failure or incident elevates to a pause. The system halts the bundle and reverts to the last known good version.
  \item \textbf{Auto-reconciliation.} Open tasks created under a paused bundle remain visible; new tasks revert to prior templates.
  \item \textbf{Receipts and lineage.} A rollback produces a receipt that links the affected users to the reverted bundle and cites the reason in plain terms.
\end{itemize}

\paragraph{Operator checklist}
\begin{enumerate}[leftmargin=*]
  \item Confirm probe or incident details on the dashboard.
  \item Hit \emph{Pause and Roll Back}. Verify that the prior version is active.
  \item Skim the auto-generated impact list and notify owners if human outreach is needed.
  \item Create a fix ticket with the failing probe attached; schedule a re-test in the digital twin.
\end{enumerate}

\subsection*{Offline Merge: Firebook/CRDT}
\paragraph{Problem}
People and teams often work with poor connectivity. Edits happen offline. When the device comes back, updates can collide.

\paragraph{Firebook}
A small, append-only log on the device records every local change with a short reason. The log is encrypted at rest and syncs when a connection returns. If sync fails partway, the device resumes from the last acknowledged entry with no data loss.

\paragraph{CRDT-based resolution}
\begin{itemize}[leftmargin=*]
  \item \textbf{Field-level merges.} Safe fields (checklist ticks, note appends) merge automatically. Structured lists use add-only sets where possible to avoid overwrites.
  \item \textbf{Priority rules.} Sensitive fields (medication dose, allergy status) never auto-merge. The system opens a side-by-side view for a human to choose, with both versions preserved until resolved.
  \item \textbf{Readable conflicts.} The reconciliation screen shows who changed what, when, and why, in plain language. Users accept, edit, or split items with one action.
  \item \textbf{Receipts.} When a conflict is resolved, the system issues a short receipt and updates lineage so auditors can replay the decision.
\end{itemize}

\paragraph{Acceptance checks}
\begin{itemize}[leftmargin=*]
  \item No local change is lost after a battery pull or app crash.
  \item Sync resumes automatically within seconds of connectivity returning.
  \item Conflict prompts appear only when needed and never block viewing critical information.
\end{itemize}

\subsection*{Operational Playbooks}
\paragraph{Webhook backlog}
\begin{itemize}[leftmargin=*]
  \item Auto-scale senders within posted limits; open a banner on dashboards showing oldest item age.
  \item If age exceeds the target window, notify partners and switch to batched delivery until the backlog clears.
\end{itemize}

\paragraph{Model hotfix}
\begin{itemize}[leftmargin=*]
  \item Freeze new releases; promote the last known good bundle.
  \item Attach failing probe snapshots to the hotfix ticket; require a twin re-run before resuming canaries.
\end{itemize}

\paragraph{Content asset failure}
\begin{itemize}[leftmargin=*]
  \item Swap to cached safe copy; mark affected lessons as \emph{degraded}.
  \item Notify content owners and restore primary assets; lift the degraded flag after a quick test.
\end{itemize}

\subsection*{Observability and Alerts}
\begin{itemize}[leftmargin=*]
  \item \textbf{Key signals.} Delivery latency, retry rate, DLQ size, rollback count, conflict rate, and time-to-resolve conflicts.
  \item \textbf{Severity.} Informational for transient retries; warning for rising backlogs; critical for data loss risk or repeated rollbacks.
  \item \textbf{User-facing notices.} When a user might notice an issue, show a short, plain banner and provide a simple workaround.
\end{itemize}

\subsection*{Success Criteria}
\begin{itemize}[leftmargin=*]
  \item High share of events delivered without manual intervention, even during partner outages.
  \item Fast, low-disruption rollbacks with clear receipts and restored stability.
  \item Offline users keep working without losing progress; conflicts resolve cleanly with minimal clicks.
  \item Incident reviews lead to specific fixes that prevent repeats.
\end{itemize}

% End Failure Modes and Compensations
% ---------- 90-Day Implementation Plan ----------
\section{90-Day Implementation Plan}

\subsection*{Overview}
This plan gets a site from decision to measurable impact in ninety days. It uses three phases: \textbf{S0} (set up and prove the loop on sandboxes), \textbf{S1} (run a small live pilot with real users), and \textbf{S2} (expand carefully, evaluate, and prepare steady-state). Each phase lists objectives, deliverables, quality checks, and exit gates. Cross-cutting streams cover instrumentation, analysis, ethics, and the runbook.

\subsection*{Phase Table}
\begin{longtable}{@{}p{0.08\linewidth}p{0.12\linewidth}p{0.28\linewidth}p{0.24\linewidth}p{0.18\linewidth}@{}}
\toprule
\textbf{Phase} & \textbf{When} & \textbf{Objectives} & \textbf{Key Deliverables} & \textbf{Exit Gate} \\
\midrule
\endhead
S0 & Weeks 1--2 & Stand up environments; connect minimum data flows; agree scope and governance; publish pilot charter. & Working sandboxes; one reference pathway configured end-to-end; localized content sample; consent and privacy notices; runbook v1. & Dry-run completes cleanly with audit entries; Steering Council approval to go live for S1. \\
S1 & Weeks 3--6 & Launch a small cohort; confirm pacing and clarity; stabilize operations; ship weekly improvements. & Live pilot; daily ops huddle; issue tracker; first dashboard; support desk scripts. & Two consecutive weeks meeting target SLOs and KPIs; no major incidents; advisory feedback logged with responses. \\
S2 & Weeks 7--12 & Expand to additional cohorts or sites; tighten targets; complete evaluation; plan steady-state. & Expansion plan; evaluation report; training packs; handoff checklist; runbook v2; public summary. & Governance sign-off on sustain/scale; budget and roles assigned; backlog triaged with owners and dates. \\
\bottomrule
\end{longtable}

\subsection*{Phase S0: Set Up and Prove the Loop}
\textbf{Objectives.} Establish a safe, minimal loop that works end-to-end on sandboxes. Agree who does what and how changes ship.\\
\textbf{Deliverables.}
\begin{itemize}[leftmargin=*]
  \item Sandbox connections for clinical and learning systems using test data.
  \item One reference pathway configured with content, tasks, and dashboards.
  \item Consent text, privacy notices, and receipts ready for pilot use.
  \item Initial SLO dashboard: availability, delivery latency, data freshness, audit completeness.
  \item Runbook draft covering incidents, on-call, and rollback steps.
\end{itemize}
\textbf{QA and checks.}
\begin{itemize}[leftmargin=*]
  \item Trace identifiers link intake, decisions, actions, and outcomes in sandbox runs.
  \item Receipts render in plain language for sample users.
  \item Equity view loads with placeholder cohorts without calculation errors.
\end{itemize}
\textbf{Exit gate.} A supervised dry-run produces expected entries with no data loss and clear explanations. Governance approves the pilot charter.

\subsection*{Phase S1: Live Pilot}
\textbf{Objectives.} Deliver value fast to a small cohort without overloading teams. Learn from real use and adjust weekly.\\
\textbf{Deliverables.}
\begin{itemize}[leftmargin=*]
  \item Live cohort enrolled with inclusion criteria and a clear opt-out route.
  \item Daily operations huddle (15 minutes) with notes and assignments.
  \item Weekly improvements: content tweaks, pacing changes, UI polish.
  \item Support desk scripts for common questions and barriers.
\end{itemize}
\textbf{QA and checks.}
\begin{itemize}[leftmargin=*]
  \item SLO tiles hold near posted targets during peak hours.
  \item Privacy pass-rate stays high with few exceptions and fast resolution.
  \item Provider acknowledgment meets the agreed window for the pathway.
  \item In-app feedback and brief calls capture user themes, logged and ranked.
\end{itemize}
\textbf{Exit gate.} Two consecutive weeks at target SLOs and KPIs, zero major incidents, and advisory feedback addressed or scheduled.

\subsection*{Phase S2: Expand and Evaluate}
\textbf{Objectives.} Grow carefully, tighten targets, and complete a practical evaluation. Prepare steady-state operations.\\
\textbf{Deliverables.}
\begin{itemize}[leftmargin=*]
  \item Additional cohorts or sites enabled with measured pacing.
  \item Evaluation report covering outcomes, equity, privacy, and operations.
  \item Training packs for clinicians, educators, and support staff.
  \item Updated runbook with lessons learned and a clear escalation tree.
  \item Public summary explaining what changed and why it matters.
\end{itemize}
\textbf{QA and checks.}
\begin{itemize}[leftmargin=*]
  \item Equity bands remain green or amber; red bands trigger mitigation before further expansion.
  \item Drift probes stable; no unexplained jumps in volume or false alerts.
  \item Receipts and audit packs export cleanly for any two-week window.
\end{itemize}
\textbf{Exit gate.} Governance approves a sustain-and-scale plan with budget, roles, and quarterly review cadence.

\subsection*{Instrumentation and Analysis}
\begin{itemize}[leftmargin=*]
  \item \textbf{Dashboards.} Post live SLO and KPI tiles with definitions and owners. Include backlog depth, retry rate, and incident counts.
  \item \textbf{Sampling.} Review a daily sample of receipts, explanations, and task outcomes for clarity and accuracy.
  \item \textbf{Equity view.} Break out engagement, adherence, and acknowledgment by approved cohorts; open mitigation tasks when gaps exceed thresholds.
  \item \textbf{Release notes.} Each weekly change includes a short note and a link to probe results that allowed release.
\end{itemize}

\subsection*{Ethics and Community Review}
\begin{itemize}[leftmargin=*]
  \item \textbf{Advisory cadence.} Meet community advisors at least twice during the pilot to preview changes and collect feedback.
  \item \textbf{User notices.} Publish plain-language notices before meaningful changes; include how to opt out or ask questions.
  \item \textbf{IRB or equivalent.} When required, maintain approvals and report summaries using generated audit packs.
\end{itemize}

\subsection*{Runbook Essentials}
\begin{itemize}[leftmargin=*]
  \item \textbf{On-call and escalation.} Define hours, roles, contact routes, and severity levels with target response times.
  \item \textbf{Incidents.} Log who, what, when, purpose, scope, impact, and fix. Send user notices when appropriate.
  \item \textbf{Rollbacks.} Keep one-click rollback ready for rules, content, or thresholds; document how to verify the prior state.
  \item \textbf{Backups and recovery.} Test backups and restore steps on a schedule; record results.
\end{itemize}

\subsection*{Weekly Cadence (Checklist)}
\begin{itemize}[leftmargin=*]
  \item Monday: review SLOs and incidents; assign fixes.
  \item Midweek: ship small improvements; update change log; verify receipts.
  \item Friday: equity review; decide next week’s focus; confirm training or outreach needs.
\end{itemize}

\subsection*{Definition of Done for the First 90 Days}
\begin{itemize}[leftmargin=*]
  \item One pathway runs end-to-end in production with posted SLOs and KPIs visible to operators and leaders.
  \item People understand and use receipts; consent changes propagate quickly.
  \item Teams acknowledge and act on tasks within expected windows without alert fatigue.
  \item Equity gaps are identified early and addressed with concrete mitigations.
  \item The runbook supports steady-state operations with a clear ownership model.
\end{itemize}

% End 90-Day Implementation Plan
% ---------- Stakeholders and Roles ----------
\section{Stakeholders and Roles}

\subsection*{Purpose}
This section names who drives which parts of the work, what decisions they own, and what “good” looks like for each group. The aim is to make collaboration simple, reduce handoff friction, and keep accountability visible.

\subsection*{Clinicians}
\textbf{Who.} Physicians, advanced practice providers, nurses, pharmacists, therapists, social workers, and care coordinators.\\
\textbf{Responsibilities.}
\begin{itemize}[leftmargin=*]
  \item Review insight cards and accept, adjust, or decline suggested actions.
  \item Keep plans current, document key decisions, and close the loop on handoffs.
  \item Watch for safety issues, escalate early, and use emergency access only when needed.
  \item Provide short feedback on clarity and burden so pathways improve over time.
\end{itemize}
\textbf{Decisions owned.}
\begin{itemize}[leftmargin=*]
  \item Whether a suggestion is appropriate now for this person.
  \item Which tasks are team-led versus self-care, and who should do them.
  \item When to pause automation for a person and when to resume.
\end{itemize}
\textbf{Tools and views.}
\begin{itemize}[leftmargin=*]
  \item A compact team console with a unified worklist, acknowledgment buttons, and quick notes.
  \item One-page summaries before visits and at discharge.
  \item Clear receipts and audit-aware updates for sensitive actions.
\end{itemize}
\textbf{Success signals.}
\begin{itemize}[leftmargin=*]
  \item High acknowledgment within expected windows without alert fatigue.
  \item Fewer duplicate asks and fewer preventable handoff errors.
  \item People report they understand “what to do next.”
\end{itemize}

\subsection*{Patients and Families}
\textbf{Who.} People receiving care, caregivers, and trusted supporters they name.\\
\textbf{Responsibilities.}
\begin{itemize}[leftmargin=*]
  \item Set preferences and goals, including language, access, and quiet hours.
  \item Use the Bring-Your-Own-Context (BYOC) pack to prepare, confirm, and follow up.
  \item Complete small steps when able, report barriers when not, and request help.
  \item Control sharing choices and review receipts in plain language.
\end{itemize}
\textbf{Decisions owned.}
\begin{itemize}[leftmargin=*]
  \item What to share, with whom, and for how long.
  \item Which channels and times work best for reminders and lessons.
  \item When to invite or remove a caregiver from shared access.
\end{itemize}
\textbf{Tools and supports.}
\begin{itemize}[leftmargin=*]
  \item A simple app or web view with checklists, short lessons, and a help button.
  \item Offline-friendly BYOC packs and one-tap language or accessibility switches.
  \item Human assistance routes for coverage, transport, devices, or cost barriers.
\end{itemize}
\textbf{Success signals.}
\begin{itemize}[leftmargin=*]
  \item High comprehension and completion for chosen actions.
  \item Fast fixes when barriers are reported.
  \item Clear understanding of data uses and easy revocation when desired.
\end{itemize}

\subsection*{Educators and Operations}
\textbf{Who.} Educators, school nurses, program leads, schedulers, support desk, and site operations managers.\\
\textbf{Responsibilities.}
\begin{itemize}[leftmargin=*]
  \item Author, localize, and version short lessons and practice activities.
  \item Schedule deliveries that respect school hours, clinic rhythms, and local calendars.
  \item Monitor engagement and comprehension, then adjust content or pacing.
  \item Triage help requests and coordinate with community partners when needed.
\end{itemize}
\textbf{Decisions owned.}
\begin{itemize}[leftmargin=*]
  \item Which cohorts to include, when to start, and when to slow or pause.
  \item Which content variants to prefer for language, literacy, and device constraints.
  \item When to open a human outreach task instead of sending another nudge.
\end{itemize}
\textbf{Tools and views.}
\begin{itemize}[leftmargin=*]
  \item An authoring studio with templates for plain language and accessibility.
  \item A learning console with role-aware dashboards and quick assignment tools.
  \item An operations panel with delivery windows, backlog, and incident banners.
\end{itemize}
\textbf{Success signals.}
\begin{itemize}[leftmargin=*]
  \item Steady engagement and comprehension without overload.
  \item Smooth schedules with few conflicts across programs.
  \item Equity gaps identified early and mitigated with concrete changes.
\end{itemize}

\subsection*{Compliance and Governance}
\textbf{Who.} Privacy officers, security leads, data stewards, quality and safety, legal, and community or ethics advisors.\\
\textbf{Responsibilities.}
\begin{itemize}[leftmargin=*]
  \item Approve purpose-bound policies and review exceptions.
  \item Oversee audits, incidents, and user notices with documented timelines.
  \item Gate higher-risk releases and verify rollbacks and receipts.
  \item Track equity measures and require mitigations when gaps appear.
\end{itemize}
\textbf{Decisions owned.}
\begin{itemize}[leftmargin=*]
  \item Go/no-go on sensitive features, data exports, and emergency overrides.
  \item Cohort definitions for equity views and minimum sizes for analysis.
  \item Acceptance of mitigation plans before resuming paused rollouts.
\end{itemize}
\textbf{Tools and views.}
\begin{itemize}[leftmargin=*]
  \item Policy console with versioned rules, purposes, and expiry windows.
  \item Audit and lineage browser with export packs for regulators or boards.
  \item Equity dashboards with banding (green/amber/red) and pause controls.
\end{itemize}
\textbf{Success signals.}
\begin{itemize}[leftmargin=*]
  \item High privacy pass-rate and complete audit trails.
  \item Few incidents and fast, well-documented resolution when they occur.
  \item Transparent change logs and regular community-facing summaries.
\end{itemize}

\subsection*{Handshakes and Expectations}
\begin{itemize}[leftmargin=*]
  \item \textbf{Clear ownership.} Each active pathway lists an accountable owner and on-call rotation.
  \item \textbf{Small, frequent changes.} Weekly improvements with short release notes and easy rollback.
  \item \textbf{Plain communication.} Short status updates, clear labels on dashboards, and simple user notices.
  \item \textbf{Feedback loops.} In-app prompts and brief meetings turn feedback into fixes within the next cycle.
\end{itemize}

% End Stakeholders and Roles


\end{document}